\documentclass[11pt,a4paper]{article}
\usepackage[T1]{fontenc}
\usepackage{lmodern}
\ifdefined\pdfgentounicode
  \input{glyphtounicode}
  \pdfglyphtounicode{parenleftbig}{0028}
  \pdfglyphtounicode{parenrightbig}{0029}
  \pdfglyphtounicode{parenleftbigg}{0028}
  \pdfglyphtounicode{parenrightbigg}{0029}
  \pdfglyphtounicode{bracketleftbigg}{005B}
  \pdfglyphtounicode{bracketrightbigg}{005D}
  \pdfglyphtounicode{vextendsingle}{007C}
  \pdfglyphtounicode{vextenddouble}{2016}
  \pdfgentounicode=1
\fi
\usepackage[margin=27mm,headheight=15pt]{geometry}
\usepackage{amsmath,amssymb,amsthm,mathtools}
\usepackage{booktabs,longtable,array,enumitem}
\usepackage{microtype,xcolor,fancyhdr}
\usepackage[colorlinks=true,linkcolor=blue!45!black,citecolor=blue!45!black,urlcolor=blue!45!black,breaklinks=true]{hyperref}
\hypersetup{pdftitle={A computer-assisted lower bound for the no-three-in-line problem},pdfauthor={},pdfsubject={A 481-phase construction attaining 1.501n for all sufficiently large n}}
\setlength{\emergencystretch}{2em}
\setlength{\parskip}{3pt}
\setlist{itemsep=2pt,topsep=4pt}
\newtheorem{theorem}{Theorem}[section]
\newtheorem{lemma}[theorem]{Lemma}
\newtheorem{proposition}[theorem]{Proposition}
\newtheorem{corollary}[theorem]{Corollary}
\theoremstyle{definition}
\newtheorem{definition}[theorem]{Definition}
\newtheorem{remark}[theorem]{Remark}
\newcommand{\E}{\mathbb E}
\newcommand{\Prob}{\mathbb P}
\newcommand{\F}{\mathcal F}
\newcommand{\R}{\mathbb R}
\newcommand{\N}{\mathbb N}
\newcommand{\Z}{\mathbb Z}
\newcommand{\ind}{\mathbf 1}
\newcommand{\dd}{\,\mathrm d}
\newcommand{\Var}{\operatorname{Var}}
\newcommand{\NoThree}{\operatorname{NoThree}}
\newcommand{\supp}{\operatorname{supp}}
\newcommand{\col}{\operatorname{col}}
\newcommand{\typ}{\operatorname{typ}}
\newcommand{\pos}[1]{\left[#1\right]_+}
\newcommand{\given}{\,\middle|\,}
\newcommand{\repo}{https://github.com/agolkjaoi/no-three-in-line-lean}
\pagestyle{fancy}
\fancyhf{}
\fancyhead[L]{\small No-three-in-line: the 481-phase construction}
\fancyhead[R]{\small 10 November 2026}
\fancyfoot[C]{\thepage}
\renewcommand{\headrulewidth}{0.3pt}
\title{\textbf{A computer-assisted lower bound\\for the no-three-in-line problem}\\[5pt]
\large A 481-phase construction attaining \(1.501n\)}
\author{}
\date{10 November 2026\\[3pt]\small Version 0.2: the 481-phase construction}
\begin{document}
\maketitle
\begin{abstract}
We give a computer-assisted proof that, for every sufficiently large integer
\(n\), the \(n\times n\) integer grid contains a set of at least
\(\lceil1.501n\rceil\) points with no three collinear. The construction uses
a finite continuous-time process whose birth rates compensate for geometric
deletion and prescribed halvings. A saved-credit martingale identifies the
expected output with deterministic input mass minus exposure lost at early
retirement. Common-stopping mixed moments and protected-line concentration
bound the five causes of this loss.

The finite input has 481 phases on a \(256\times256\) spatial partition,
four parity types, and exact mass \(1.501114506104296875\).
Its capacity and root-deadline risk bounds are checked using integer and
rational upper bounds over entire spatial cells and angular intervals.
The accompanying package contains the input, checker source and full fresh
replay records. Fixed moment orders and a total loss budget of \(10^{-5}\)
suffice for the final comparison. No divisibility restriction on \(n\) is
used. The argument is asymptotic and does not provide a practical numerical
threshold for \(n\).
\end{abstract}
\noindent\textbf{Reading and attribution.}
No knowledge of Lean is needed for the mathematical text.
The probabilistic arguments are proved in ordinary mathematical language;
the executable checks concern the finite geometric input.
AI systems contributed substantially to the research, code, proofs, checking
and exposition. No named author is supplied. This manuscript has not received
independent human peer review, and the separate partial Lean development is
not a formal proof of its final theorem. Precise disclosure and execution
scope appear in Section~\ref{sec:status}.
\tableofcontents
\newpage
\section{The target and the structure of the argument}\label{sec:overview}

Write \(V_n=\{0,\ldots,n-1\}^2\). A set \(S\subseteq V_n\) has the
\emph{no-three-in-line property} if every three distinct points of \(S\)
are noncollinear. Equivalently,
\[
 (q_x-p_x)(r_y-p_y)-(q_y-p_y)(r_x-p_x)\ne0
 \qquad(p,q,r\in S\text{ distinct}).
\]
Every such set has at most two points on each horizontal row, hence
\(|S|\le2n\). The numerical conclusion proved below is
\begin{equation}\label{eq:target}
 \exists n_0\ \forall n\ge n_0\ \exists S\subseteq V_n:
 \quad \NoThree(S),\qquad 1501n\le1000|S|.
\end{equation}
The positive improvement over \(3n/2\) is the fixed number \(n/1000\).
The quantifier is over every sufficiently large integer \(n\), with no
divisibility restriction.

\subsection{Finite input, fixed before the grid size}
Choose a finite rational time partition
\(0=t_0<t_1<\cdots<t_J=T\) and a fixed finite rational rectangular
partition of the unit square. A common refinement may be taken to be a
\(G\times G\) grid for a fixed integer \(G\). Let \(\mathcal B\) be the
union of its horizontal and vertical grid lines.
There are four parity types \(r\in(\Z/2\Z)^2\). For each type choose
a deterministic function
\[
 f_{r,t}:[0,1]^2\longrightarrow[0,1]
\]
that is constant on every open spatial cell and every time interval.
Set it to zero on \(\mathcal B\), outside the square, and after time \(T\).
Define
\begin{equation}\label{eq:input}
 m_t(p)=f_{\typ(p),t}(p/n),\qquad
 F_{r,t}(x)=\int_0^t f_{r,u}(x)\dd u,\qquad
 A=\frac14\sum_r\int_{[0,1]^2}F_{r,T}(x)\dd x.
\end{equation}
Here \(\typ(p)=p\bmod2\). The factor \(1/4\) belongs to the spatial
mass, not to the birth rate of an individual point.
The last-use time \(\tau_r(x)\) is the right endpoint of the last interval
on which that type and cell is used; it is zero for an unused cell.
For a grid point write \(\tau_p=\tau_{\typ(p)}(p/n)\).

The concrete construction has \(G=256\) and 481 phases. An integer
field and a rank assigned to each positive type-cell specify its one
exposure phase. Ranks are used in decreasing order, so the cumulative
field at the end of the phase of rank \(j\) is supported on ranks
at least \(j\). Section~\ref{sec:certificates} gives the exact masks,
durations and data format. Removing \(\mathcal B\) is part of the
definition of the discrete process, even though it has zero planar area.

\subsection{The geometric input conditions}
A primitive integer direction \(v\) is identified with \(-v\). Put
\(h(v)=\max(|v_1|,|v_2|)\) and \(\bar v=v\bmod2\ne0\).
For a field \(F\) define the major-coordinate chord integral
\[
 C_vF(x)=\int_{\R}F\bigl(x+s v/h(v)\bigr)\dd s.
\]
This is not an arc-length integral. Let \(\Gamma\) be a fixed finite
family of designated directions containing both coordinate directions,
of cardinality \(g\).
The numerical example uses all directions of height at most three,
so \(g=16\). Define the line capacity
\begin{equation}\label{eq:capacity}
 L_{v,r}(F;x)=\frac{C_v(F_r+F_{r+\bar v})(x)}{2h(v)}.
\end{equation}
Let the four direction families be
\((1,u),(1,-u),(u,1),(-u,1)\), \(0\le u\le1\), and let
\(c_{F,j}(x,u)\) denote the corresponding chord. Set
\begin{equation}\label{eq:risk}
 g_F(x)=\frac14\sum_{j=1}^4\int_0^1c_{F,j}(x,u)^2\dd u,
 \qquad
 \mathcal R_r(F;x)=\frac{D_*}{6}\sum_{s\ne r}g_{F_r+F_s}(x).
\end{equation}
Fix a finite set \(P\) of pairwise coprime odd integers greater than one
and require \(D_*\ge D_P=\tfrac34\prod_{q\in P}(1-q^{-2})\).
This is part of the admissibility hypothesis below.
Section~\ref{sec:sieve} proves the resulting finite-sieve bound;
the concrete input uses \(D_*=76/125\) and 175 odd prime factors.

The finite schedule is \emph{admissible} for fixed parameters if there
are positive constants \(\kappa_H,\kappa_L\) and a number \(0\le r_*<d\)
such that
\begin{align}
 L_{v,r}(F_T;x)&\le K-\kappa_H &&(v\in\Gamma),\label{eq:heavyinput}\\
 L_{v,r}(F_T;x)&\le\lambda-\kappa_L &&(v\notin\Gamma),\label{eq:lightinput}\\
 \mathcal R_r(F_{\tau_r(x)};x)&\le r_* &&(\tau_r(x)>0),\label{eq:riskinput}
\end{align}
uniformly over the relevant roots, types and lines. In the last formula
\emph{every one of the four fields is evaluated at the same root's
deadline}. The constants satisfy
\begin{equation}\label{eq:exponents}
 \begin{gathered}
 0<K<2,\quad 0<\rho<1,\quad 0\le\lambda<1-\rho,\\
 0<d<\gamma<1,\quad 0<\xi<d,\quad \zeta>0,\quad 1-d-2\xi-\zeta>0.
 \end{gathered}
\end{equation}
All data and all displayed constants are fixed before \(n\) tends to
infinity. Constants depending on a moment order may depend on that
order, but the order is also fixed before taking the limit.

\begin{theorem}[Analytic reduction]\label{thm:bridge}
For an admissible finite schedule, the stopped compensated process
defined in Section~\ref{sec:process} produces a random set \(S_n\subseteq
V_n\) with the no-three-in-line property and
\[
 \E|S_n|=(A-o(1))n.
\]
Consequently, for every fixed \(c<A\), a set of size at least \(cn\)
exists for every sufficiently large \(n\).
\end{theorem}

The detailed argument for this theorem occupies
Sections~\ref{sec:geometry}--\ref{sec:closing}. The concrete schedule
and the exact finite verification of
\eqref{eq:heavyinput}--\eqref{eq:riskinput} are given in
Sections~\ref{sec:certificates} and~\ref{sec:numerics}.
Corollary~\ref{cor:target} combines them with the analytic reduction
to obtain \eqref{eq:target}. The finite computation verifies the
geometric input; the probabilistic lemmas are proved in the text.

\subsection{Why saved credit is the useful quantity}
Each active point has a weight \(w_p\), initially one. It is selected at
rate \(m_t(p)w_p/n\). Weights grow to compensate for the chance of
selection, deletion by a newly completed pair, or a prescribed
halving. At an early safety stop, the positive weight is saved rather
than discarded. The resulting process \(Z_p\) is a bounded martingale
with \(\E Z_p(t)=1\).

The change of probability measure
\[
 \dd\Prob^{(p)}=Z_p(T)\dd\Prob
\]
measures the amount of credit lost by retiring \(p\), rather than just
the raw probability that \(p\) retires. Let \(\theta_p\) be its first
positive saving time and \(E_p\) the event that this happens strictly
before \(\tau_p\) for a safety reason. Put
\[
 A_n=\frac1{n^2}\sum_{p\in V_n}F_{\typ(p),T}(p/n).
\]
The exact identity proved in Section~\ref{sec:loss} is
\begin{equation}\label{eq:previewloss}
 A_n-\frac{\E|S_n|}{n}
 =\frac1{n^2}\sum_p\E^{(p)}\left[
 \ind_{E_p}\int_{\theta_p}^{T}m_t(p)\dd t\right].
\end{equation}
Thus a bound \(\sup_p\Prob^{(p)}(E_p)\le\varepsilon\) implies
\(\E|S_n|/n\ge(1-\varepsilon)A_n\).
There are five causes to control: weight cap, ordinary-line tag,
pair-correction floor, virtual-deletion bound, and designated-line clock.
The last three are global safety stops. The paper proves their estimates
on the already stopped process; it does not assume that those stops
are rare while deriving the estimates that make them rare.

\subsection{A fixed budget suffices for the numerical target}
The exact mass and target are
\begin{equation}\label{eq:massnumbers}
 A =\frac{3842853135627}{2560000000000},\qquad
 c=\frac{1501}{1000},\qquad
 \Delta=A-c=\frac{293135627}{2560000000000}.
\end{equation}
The grid comparison gives \(A_n\ge A(1-512/n)^2\). For sufficiently
large \(n\), the five stopping-cause estimates with fixed parameters give
\(\sup_p\Prob^{(p)}(E_p)<10^{-5}\). Taking also \(n\ge2^{24}\),
\[
 (1-10^{-5})A(1-2^{-15})^2-c
 =\frac{2165079963272809251797}{274877906944000000000000000}>0.
\]
This strict rational inequality is enough for \eqref{eq:target}.
There is no need to prove an effective, practically small threshold
\(n_0\), or to compute an unbounded sequence of ever tighter error
budgets. The asymptotic argument is a way to obtain one fixed budget.
The condition \(n\ge2^{24}\) bounds only the displayed mass error;
the probability estimates impose additional asymptotic thresholds.


\section{Uniform lattice geometry}\label{sec:geometry}

In this section a line means its intersection with \(V_n\), and only
lines containing at least two grid points are counted. Write \(L_\ell\)
for its cardinality and \(L_{pq}\) for the length of the line through
distinct \(p,q\). Constants in this section are uniform in the root's
position, including roots approaching a cell boundary as \(n\) changes.

\begin{lemma}[Direction and line counts]\label{lem:lines}
There are at most \(4h\) primitive unoriented directions of height
\(h\). Every grid line of that height has at most \(2n/h\) points.
Consequently, uniformly in \(p\),
\begin{align}
 \sum_{\ell\ni p}L_\ell^2&\le16n^2(1+\log n),\label{eq:line2}\\
 \sum_{\ell\ni p}L_\ell^k&\le
 4(2n)^k\left(1+\frac1{k-2}\right)\qquad(k>2).
 \label{eq:linek}
\end{align}
\end{lemma}
\begin{proof}
If \(v\) is primitive, the integer points on its line are precisely
\(p+\Z v\). Indeed, Bezout coefficients show that a real multiple of
\(v\) with integer coordinates has an integer multiplier. A coordinate
of absolute step \(h\) gives \(L_\ell\le1+(n-1)/h\le2n/h\).
For \(h\ge2\), counting the primitive points on the boundary of the
square of radius \(h\), with opposite signs identified, gives
\(4\varphi(h)\le4h\); height one has four directions. Sum
\(4h(2n/h)^k\), using the harmonic bound when \(k=2\) and an integral
comparison when \(k>2\). The same argument gives
\(\sum_{q\ne p}L_{pq}=\sum_{\ell\ni p}(L_\ell-1)L_\ell
\le\sum L_\ell^2\).
\end{proof}

\begin{lemma}[Two uniform incidence kernels]\label{lem:kernels}
For every fixed \(\xi>0\),
\begin{align}
 \sum_{q\notin\ell(a,b)}L_{aq}L_{bq}&\le C_\xi n^{2+\xi}
       &&(a\ne b),\label{eq:twostar}\\
 \sum_{q\in\ell\cap V_n}L_{pq}&\le C_\xi n^{1+\xi}
       &&(p\notin\ell).
       \label{eq:fixedline}
\end{align}
\end{lemma}
\begin{proof}
The elementary divisor estimate \(\tau(m)\le C_\epsilon m^\epsilon\)
follows by writing \(m=\prod l^{a_l}\). For sufficiently large primes,
\(a+1\le2^a\le l^{\epsilon a}\); the finitely many smaller primes have
finite suprema of \((a+1)/l^{\epsilon a}\). Hence, with
\(\Delta_n=\max_{m\le2n^2}\tau(m)\),
\(\Delta_n(1+\log n)^2\le C_\xi n^\xi\).

For \eqref{eq:twostar}, write \(q-a=e u\), \(q-b=f v\), where
\(e,f>0\) and \(u,v\) are primitive, and put \(w=b-a\).
Then \(eu-fv=w\) and \(f\mid\det(u,w)\). The exclusion of the baseline
ensures \(1\le|\det(u,w)|\le2n^2\), so for fixed \(u\) there are at
most \(\Delta_n\) choices of \(f\). In dyadic height bands
\(H\le h(u)<2H\), \(J\le h(v)<2J\), assume \(H\le J\).
There are \(O(H^2)\) choices of \(u\). The congruence
\(eu\equiv w\pmod f\) confines \(e\) to at most one residue class
modulo \(f\), by Bezout. A coordinate with \(|u_i|\ge H\), together
with \(|v_i|<2J\), confines \(e/f\) to an interval of length at most
\(4J/H\). Thus the number of candidates is
\(O(\Delta_n HJ)\). Their individual weights are at most
\(4n^2/(HJ)\). Interchange \(a,b\) when \(H>J\), then sum the
\(O((1+\log n)^2)\) bands.

For \eqref{eq:fixedline}, the cases of zero or one grid point are
immediate. Otherwise write \(\ell\cap\Z^2=a+\Z v\) with primitive
\(v\), and let \(D=|\det(v,p-a)|\in[1,2n^2]\).
For \(q\in\ell\), put \(g(q)=\gcd(|q_1-p_1|,|q_2-p_2|)\).
Then \(g(q)\mid D\). For each \(e\mid D\), the condition
\(q\equiv p\pmod e\) restricts the line parameter to one residue class.
Within \(\|q-p\|_\infty\le R\), this gives at most
\(2R/(e h(v))+1\) choices. Since
\(g(q)\le\sum_{e\mid g(q)}e\) and relevant \(e\le R\),
\[
 \sum_{\|q-p\|_\infty\le R}g(q)\le3R\tau(D).
\]
In the shell \(R/2<\|q-p\|_\infty\le R\),
\(L_{pq}\le4ng(q)/R\), so the shell contributes at most
\(12n\tau(D)\). Sum dyadic shells and use the divisor estimate.
\end{proof}

\subsection{Parity, chords and a uniform boundary estimate}
Suppose \(F_r\) are constant on a fixed partition into \(R\) open
rectangular cells, vanish on the common boundaries, and lie in
\([0,M]\). Along a primitive line through a root of type \(r\), the
types alternate between \(r\) and \(r+\bar v\).

\begin{lemma}[Lattice-to-chord comparison]\label{lem:chord}
For every root and primitive direction of height \(h\),
\begin{equation}\label{eq:chorderror}
 \left|\sum_{q\in\ell(p,v)\setminus\{p\}}
 F_{\typ(q)}(q/n)
 -\frac n{2h}C_v(F_r+F_{r+\bar v})(p/n)\right|
 \le(4R+1)M.
\end{equation}
Moreover \(A_n=A+O_{R,T}(n^{-1})\).
\end{lemma}
\begin{proof}
Each parity subsequence has major-coordinate spacing \(2h/n\).
The number of points of any shifted one-dimensional lattice in an
open interval differs from interval length divided by spacing by at
most two. Sum over the two types and the \(R\) cells, then remove the
root, costing at most \(M\). A line coinciding with a common boundary
has field identically zero on both sides of \eqref{eq:chorderror}.
For the mass statement, count each coordinate of each parity class
in an open rectangle: the count is
\(n^2\operatorname{area}(E)/4+O(n)\). Sum the fixed number of cells
and types. No assumption that \(n\) is even or a multiple of \(G\)
is used.
\end{proof}

The regular grid in the concrete input permits a useful relative
version of the mass estimate.
\begin{lemma}[Mass on every integer grid]\label{lem:massgrid}
For nonnegative four-type fields constant on the open cells of a
\(G\times G\) grid and zero on its grid lines, let
\(A=\tfrac14\sum_r\int F_r\) and
\(A_n=n^{-2}\sum_{p\in V_n}F_{\typ(p)}(p/n)\). For every integer
\(n\ge2G\),
\begin{equation}\label{eq:massgrid}
 A(1-2G/n)^2\le A_n\le A(1+2G/n)^2.
\end{equation}
\end{lemma}
\begin{proof}
For a coordinate parity \(e\in\{0,1\}\), the number of integers of
that parity in \((in/G,(i+1)n/G)\) is
\[
 N_{i,e}=\left\lceil\frac{(i+1)n-eG}{2G}\right\rceil
          -\left\lfloor\frac{in-eG}{2G}\right\rfloor-1.
\]
It lies between \(n/(2G)-1\) and \(n/(2G)+1\). Since the lower
bound is nonnegative, the number of points of any specified parity
type in a cell lies between the squares of these bounds. Multiply
by that cell's nonnegative field value and sum over cells and types.
The continuum contribution of a typed cell of value \(v\) is
\(v/(4G^2)\), giving \eqref{eq:massgrid}. Boundary points have field
zero, as required by the definition of the input.
\end{proof}

\begin{lemma}[Uniform bounded variation]\label{lem:bv}
For a nonnegative field \(F\le M\) on \(R\) fixed rectangular cells,
each of the four functions \(u\mapsto c_{F,j}(x,u)^2\) satisfies
\[
 \|c_{F,j}(x,\cdot)^2\|_\infty\le M^2,
 \qquad \Var(c_{F,j}(x,\cdot)^2)\le(4R+2)M^2,
\]
uniformly in \(x\in[0,1]^2\).
\end{lemma}
\begin{proof}
For \((1,u)\), the chord of a rectangle \((a,b)\times(c,d)\) through
\((x,y)\), for \(u>0\), is
\[
 \pos{\min(b,x+(d-y)/u)-\max(a,x+(c-y)/u)}.
\]
Its changes of branch occur only at slopes from the root to the four
vertices. Commonly refine the at most \(4R\) switch values. On each
interval the total chord has the form \(\alpha+\beta/u\), hence is
monotone, as is its square since it is nonnegative. The square lies
in \([0,M^2]\). At positive switches the values agree. At zero, a line
may coincide with a common boundary; its possible endpoint jump costs
at most one further \(M^2\). The stated bound follows. Reflections
and coordinate exchange give the other three families.
\end{proof}

This proof uses bounded variation, not a uniform derivative bound.
For example, for a half-square and a root a distance \(\epsilon\)
above its boundary, a chord can have the form
\(\max(0,1/2-\epsilon/u)\). Its derivative near \(u=2\epsilon\)
is unbounded as \(\epsilon\downarrow0\), whereas its squared
variation remains at most \(1/4\).

\subsection{A finite sieve for primitive directions}\label{sec:sieve}
Let \(P\) be a fixed finite set of pairwise coprime odd integers greater
than one; primality is unnecessary. Put
\[
 Q=\prod_{p\in P}p,\qquad
 D_P=\frac34\prod_{p\in P}(1-p^{-2}),\qquad
 A_P=\prod_{p\in P}(1+p^{-2}).
\]
For \(e=(\epsilon,\eta)\in\{(0,1),(1,0),(1,1)\}\) and a nonnegative
BV function \(w\) on \([0,1]\), define
\[
 K_{N,e}(w)=\sum_{h\le N}\frac1{h^2}
 \sum_{\substack{1\le a\le h\,;\ (a,h)=1\\(h,a)\bmod2=e}}w(a/h).
\]
Write \(I=\int_0^1w\), \(M_w=\sup w\), and \(V_w=\Var(w)\).

\begin{lemma}[Finite-sieve upper bound]\label{lem:sieve}
For every real \(N\ge Q\),
\begin{equation}\label{eq:sieve}
 K_{N,e}(w)\le\frac{D_P}{3}I\log N+C_P(M_w+V_w),\qquad
 C_P=\left(5+\frac14\log Q\right)A_P.
\end{equation}
The rational upper bound
\(C_P\le[5+\tfrac14\sum_{p\in P}(p-1)]A_P\) is also valid.
\end{lemma}
\begin{proof}
For real \(x\ge1\), omit the coprimality condition and denote the
corresponding sum by \(J_e(x;w)\). Counting alternating tags gives
\[
 \left|\sum_{\substack{1\le b\le k\\b\equiv\eta\ (2)}}w(b/k)
       -\frac k2 I\right|\le2M_w+V_w.
\]
To see this, place intervals of length \(2/k\) to the left of the
tags; the oscillation errors sum to at most \(V_w\), and the omitted
or cut end intervals have total length at most \(4/k\).
Also the parity harmonic sum differs from \(\tfrac12\log x\)
by at most two, by comparing ordinary harmonic sums with integrals.
The finite estimate \(\sum_{k\le m}k^{-2}\le2\) therefore gives
\[
 J_e(x;w)=\frac I4\log x+E_e(x),\qquad
 |E_e(x)|\le5(M_w+V_w).
\]
Nonnegativity of \(w\) allows the bound
\[
 \ind_{(a,h)=1}\le
 \prod_{p\in P}(1-\ind_{p\mid a,\ p\mid h}).
\]
Expand this finite product. For \(A\subseteq P\), put
\(d_A=\prod_{p\in A}p\) and \(s_A=(-1)^{|A|}\).
Pairwise coprimality identifies simultaneous divisibility with
divisibility by \(d_A\); oddness preserves the parity class. Thus
\[
 K_{N,e}(w)\le\sum_{A\subseteq P}\frac{s_A}{d_A^2}J_e(N/d_A;w).
\]
Since \(N\ge Q\), every \(N/d_A\ge1\). Substituting the two-sided
estimate for \(J_e\) and using
\(\sum s_A/d_A^2=4D_P/3\), \(\sum1/d_A^2=A_P\),
\(\log d_A\le\log Q\), and \(I\le M_w\) proves
\eqref{eq:sieve}. It is essential to use a two-sided error estimate
inside the signed inclusion-exclusion sum. Finally \(\log p\le p-1\)
gives the rational constant.
\end{proof}

Only a one-sided bound is needed in this application. Consequently
no infinite Mobius series, value of \(\zeta(2)\), or assertion of
probabilistic independence between parity and direction is required.
The recorded choice of 175 factors gives \(D_P<76/125\); its product
\(Q\) has 440 decimal digits. This large fixed threshold is compatible
with the asymptotic quantifiers, but gives no useful small-grid claim.
The factors and rational comparison are included in the companion data.

\subsection{The common-deadline pair sum}
For four fields \(F\) and a root \(p\), define
\begin{equation}\label{eq:pairsum}
 \mathcal B_n(p;F)=\frac1{2n^2}
 \sum_{\substack{q,z\in V_n\setminus\{p\},\ q\ne z\\p,q,z\text{ collinear}}}
 F_{\typ(q)}(q/n)F_{\typ(z)}(z/n).
\end{equation}
The pair sum is ordered; the factor \(1/2\) removes this duplication.
A superscript \(\Gamma^c\) restricts it to nondesignated directions.

\begin{proposition}[Uniform risk comparison]\label{prop:pair}
If \(D_*\ge D_P\), then for \(n\ge\max(2,Q)\),
\[
 \mathcal B_n^{\Gamma^c}(p;F)
 \le\mathcal R_{\typ(p)}(F;p/n)\log n+C_{\rm geom},
\]
where, with \(E_0=(4R+1)M\), one may take
\begin{equation}\label{eq:geomconstant}
 C_{\rm geom}=M^2+4ME_0+2E_0^2+6C_P(4R+3)M^2.
\end{equation}
This constant is independent of \(n,p\) and of the field values
within their fixed partition and bound.
\end{proposition}
\begin{proof}
For each direction let \(W_v\) be the weighted line sum with the root
removed. The contribution is
\((W_v^2-\sum_qF_{\typ(q)}(q/n)^2)/(2n^2)\).
Discard the nonnegative diagonal term. By \eqref{eq:chorderror},
\(W_v\le nC_v(F_r+F_{r+\bar v})/(2h)+E_0\), while
\(C_v(F_r+F_{r+\bar v})\le2M\). Expanding the square, and using
at most \(4h\) directions at height \(h\le n\), gives
\[
 \mathcal B_n^{\Gamma^c}(p;F)
 \le\frac18\sum_{\substack{v\text{ primitive}/\pm\\h(v)\le n}}
 \frac{C_v(F_r+F_{r+\bar v})(p/n)^2}{h(v)^2}
       +4ME_0+2E_0^2.
\]
Indeed, the cross term is at most
\(\sum_{h\le n}4h\,ME_0/(nh)\le4ME_0\); the squared error is at
most \(\sum_{h\le n}4h\,E_0^2/(2n^2)\le2E_0^2\).
Represent nonaxis directions by the four families with slope
\(a/h\), \(1\le a\le h\). The diagonal overlap can be retained for
an upper bound; the two axes contribute at most \(M^2\) separately.
Apply Lemma~\ref{lem:sieve} to each of the twelve family/parity
functions. Their supremum plus variation is at most
\((4R+3)(2M)^2\), by Lemma~\ref{lem:bv}, so the total sieve error,
including the factor \(1/8\), is at most \(6C_P(4R+3)M^2\).
For each nonzero parity \(v\), the four integrals sum to
\(4g_{F_r+F_{r+v}}\). The coefficient is
\(\tfrac18\cdot\tfrac{D_P}{3}\cdot4=D_P/6\).
Increase it to \(D_*/6\). The directions discarded at the start
only increase the upper bound.
\end{proof}

For a general time-dependent input take \(M=T\) once and for all.
Any smaller bound on its terminal fields is also valid. The same
constant works for every time, including every root's last-use time.
For a fixed root with deadline \(\tau\), finite-product differentiation
gives the exact identity
\begin{equation}\label{eq:timesym}
 \frac1{n^2}\int_0^\tau\sum_{q,z}^{\Gamma^c}m_t(q)F_{\typ(z),t}(z/n)\dd t
 =\mathcal B_n^{\Gamma^c}(p;F_\tau).
\end{equation}
Indeed, differentiate each product \(F_t(q)F_t(z)\), exchange \(q,z\),
and integrate. The index set is symmetric and consists of distinct
points different from the root. Using separate deadlines for \(q\)
and \(z\) would not give this identity.

\subsection{Explicit constants for the 481-phase input}\label{sec:uniform481}
The input defined in Section~\ref{sec:rankinput} has \(G=256\).
Every cumulative field is dominated by its terminal field, whose
maximum is \(6.25798212168<7\). Thus \(R=256^2=65536\) and \(M=7\)
work at every time, although the sum of the phase durations is much
larger than seven. Lemma~\ref{lem:massgrid} gives
\[
 A_n\ge A(1-512/n)^2\qquad(n\ge512).
\]

For the 175 odd primes from 3 through 1049 used by the finite sieve,
\(\sum_{p\in P}(p-1)=84164\). Also
\[
 S:=\sum_{p\in P}p^{-2}<\sum_{k\ge1}\frac1{4k(k+1)}=\frac14.
\]
The elementary symmetric sum of degree \(j\) in the numbers \(p^{-2}\)
is at most \(S^j\); expanding the finite product gives
\(A_P\le(1-S)^{-1}<4/3\). Hence the rational choice
\(C_P=84184/3\) is sufficient in Lemma~\ref{lem:sieve}. Substitution
in \eqref{eq:geomconstant}, with \(E_0=1835015\), gives
\begin{equation}\label{eq:geomnumeric}
 C_{\rm geom}=8897332619623<10^{13}.
\end{equation}

Consequently a continuum certificate
\(\mathcal R_{\typ(p)}(F_{\tau_p};p/n)\le r_*=9999/10000\)
implies, uniformly over every used root,
\begin{equation}\label{eq:discrete481}
 \mathcal B_n^{\Gamma^c}(p;F_{\tau_p})\le a\log n,
 \qquad a=\frac{9999}{10000}+\frac1{320000}
          =\frac{319969}{320000}<\frac{19999}{20000}=d,
\end{equation}
provided \(\log n\ge3\cdot10^{18}\). This condition implies
\(n\ge Q\), since \(\log Q\le84164\), and it absorbs
\eqref{eq:geomnumeric} with room to spare. It is only a sufficient
threshold for the deterministic lattice comparison, not a claimed
threshold for all stochastic error terms in the final theorem.


\section{The stopped compensated process}\label{sec:process}

Fix \(n\ge2\) and put \(B=n^d\), \(M_n=n^\gamma\), and
\(\epsilon_n=n^{-\zeta}\). A point is active, selected, geometrically
dead, or positively retired. Initially every point is active, its weight
is one, and \(S=\varnothing\). Selection is permanent. Dead and retired
points never return to the active set \(Q_t\).

\subsection{Line variables, rates and flows}
Every designated line has occupancy \(N_\ell=|S\cap\ell|\) and clock
\(k_\ell\). For \(N=0,1\) set
\[
 a_N(k)=\frac{2-N}{2-k},\qquad b_0=\frac12,\quad b_1=1.
\]
Every nondesignated line has a credit \(r_\ell\), initially one.
It changes only while the line contains exactly one selected point.
Before the global stopping time \(\sigma\), define
\begin{align}
 \lambda_p(t)&=\ind_{p\in Q_t}\frac{m_t(p)w_p(t)}n,
 &\Lambda_\ell&=\sum_{p\in\ell}\lambda_p,\label{eq:rates}\\
 \dot k_\ell&=\Lambda_\ell/a_{N_\ell}(k_\ell)
       &&(N_\ell=0,1),\label{eq:clockflow}\\
 \dot r_\ell&=-r_\ell\Lambda_\ell
       &&(N_\ell=1,\ \ell\text{ nondesignated}).\label{eq:tagflow}
\end{align}
The other clock and line-credit derivatives are zero. For active \(p\),
put
\[
 H_p=\sum_{q\in Q_t\setminus\{p\}}\lambda_q
 \ind_{\ell(p,q)\notin\Gamma,\ |S\cap\ell(p,q)|=1}.
\]
Here \(\ell\notin\Gamma\) means that its direction is nondesignated.
The weight flow is
\begin{equation}\label{eq:weightflow}
 \frac{\dot w_p}{w_p}=H_p+\lambda_p+
 \sum_{\substack{\ell\ni p\\\ell\in\Gamma}}
 (\Lambda_\ell-\lambda_p)b_{N_\ell}.
\end{equation}
The term \(H_p\) counts each actual cause \(q\) once, not once for
each possible witness. All rates are set to zero after \(\sigma\).

At a birth of \(q\), use the pre-birth state to choose the branch,
then perform these operations in order:
\begin{enumerate}
\item Select \(q\) and set its weight to zero.
\item Kill every active point that now lies on a line with two selected
points, setting its weight to zero.
\item On each designated line having its first birth, halve the weights
of the remaining active points.
\item Update occupancies and the activation flags of ordinary lines.
The values of \(k_\ell\) and \(r_\ell\) are not reset.
\item Process any newly reached safety boundary and positive retirement.
\end{enumerate}
For a fixed other point \(p\), at most one line through \(p,q\) can
cause a halving. Geometric death takes priority over positive saving.

\subsection{Safety boundaries and the observation convention}
A point retires positively at its first weight-cap hit \(w_p=B\),
at a hit \(r_\ell=\rho\) on an ordinary line through it, or at its
deadline \(\tau_p\). Its current positive weight is saved. Unused
points retire at time zero with saved weight one.

Two additional variables define global guards. For distinct active
\(a,b\), let \(\kappa_{ab}\) be the sum of rates of birth causes
\(q\notin\ell(a,b)\) that would change both weights by a factor less
than one. Set \(\dot\eta_{ab}=-\kappa_{ab}\eta_{ab}\), initially
\(\eta_{ab}=1\), and freeze this variable when either endpoint leaves
the active state. For a designated line \(\ell\) and an active point
\(q\notin\ell\), define the virtual deletion mass
\begin{equation}\label{eq:virtual}
 \mathcal V_{\ell,q}=
 \sum_{y\in\ell\cap Q_t}\frac{w_y}{a_{N_\ell}(k_\ell)}
 \ind_{\ell(q,y)\notin\Gamma,\ |S\cap\ell(q,y)|=1}.
\end{equation}
Globally stop at the first hit of \(k_\ell=K\),
\(\eta_{ab}=1-\epsilon_n\), or \(\mathcal V_{\ell,q}=M_n\), including
an exceedance produced by a jump. Save the positive active credits
and freeze the process. If no guard is hit, \(\sigma=\infty\);
the deterministic terminal time \(T\) is not itself a failure event.

When a mixed observation is saved at a root's positive boundary, take
its value after geometric updates but before removing auxiliary roots
because of simultaneous positive retirement. At a continuous hit take
the boundary values before changing the flags. This convention prevents
a tag witness from disappearing at the very instant it is observed.
It does not alter the selected set or its birth rates. Auxiliary roots
that retired strictly earlier contribute zero.

\begin{proposition}[Construction and legality]\label{prop:construction}
For fixed \(n\), the stopped process is well-defined on \([0,T]\),
is adapted with right-continuous paths and left limits, and always
produces a no-three-in-line set of at most \(2n\) points.
\end{proposition}
\begin{proof}
Give each point a Poisson candidate clock of rate \(B/n\), together
with independent uniform marks. At a candidate arrival, accept the
birth if its mark is at most \(m_t(p)w_p(t-)/B\) and the point is
eligible in the pre-arrival state. The total candidate rate is \(nB\),
so there are finitely many arrivals on \([0,T]\) almost surely.
Use the filtration revealing only arrivals and marks up to the current
time, with the usual completion and right-continuous modification.

Between arrivals and mode changes, solve the finite-dimensional ODE.
Within the stopped region the denominators are separated from zero:
\(k\le K<2\), \(r\ge\rho>0\), and \(\eta\ge1-\epsilon_n>0\).
Using \(w\), rather than \(\log w\), as a coordinate, each mode has a
locally Lipschitz vector field. Extend it outside a compact
neighborhood with a cutoff if necessary. The resulting flow is unique
and depends continuously on its initial state. First hits of its
closed guards are measurable: on any bounded time interval, the
minimum distance of the continuous flow to a closed guard detects a
hit. Composition at the preceding stopping time preserves adaptation.

Every positive retirement occurs at most once per point. Occupancies
and line activation flags change only finitely often, deadlines and
phases are finite, and a global stop occurs once. Resolve all
simultaneous guards together and mark them as processed. Together with
the finite candidate arrivals, this rules out an accumulation of
updates. The construction gives the claimed path regularity.

Immediately before every birth the chosen point is active, hence is
not on a line with two selected points. Adding it preserves legality;
the subsequent deaths maintain this invariant. The row bound gives
\(|S|\le2n\).
\end{proof}

For compensators, evaluate the intensity in the pre-jump state and use
\(\ind_{t\le\sigma}\). The birth that triggers a global stop is counted.
In time integrals one may instead use \(\ind_{t<\sigma}\), since the
two differ on at most one time per path. Thinning gives, for a
nonnegative predictable observation \(g\),
\begin{equation}\label{eq:compensator}
 \E\sum_{\text{births }(p,t)}g_p(t)
 =\E\int_0^T\sum_p\lambda_p(t)g_p(t)\dd t.
\end{equation}
For simple predictable functions this is independent increments and
the uniform-mark rule; monotone approximation gives the general
identity. Bounded signed observations follow by subtraction. All
fixed-\(n\) rates and stopped observations used below are integrable.

\section{Saved credit and the exact output identity}\label{sec:loss}

Define \(Z_p\) to equal \(w_p\) while \(p\) is active, zero after
selection or geometric death, and its saved value after positive
retirement or a global stop. Let \(C_p\) be the birth-eligible credit,
including the global-stop indicator, and let \(V_p\) be saved credit.
Then \(Z_p=C_p+V_p\) and \(0\le Z_p\le B\).

\begin{lemma}[One-point compensation]\label{lem:credit}
For each point \(p\), \(Z_p\) is a bounded martingale and
\(\E Z_p(t)=1\).
\end{lemma}
\begin{proof}
While \(p\) is active, let \(J_p(q)\in\{0,1/2,1\}\) be its
multiplier at a birth of \(q\). Selection or deletion gives zero,
a first designated-line birth gives \(1/2\), and otherwise it is one.
The flow \eqref{eq:weightflow} is precisely
\[
 \dot w_p=w_p\sum_q\lambda_q(1-J_p(q)).
\]
By \eqref{eq:compensator}, the jump drift is its negative. Positive
saving changes no value; subsequent drift is zero. A coincident
geometric death has already been assigned multiplier zero. Thus the
compensated drift vanishes, and boundedness makes the result a true
martingale. Its initial value, including immediate saving of unused
points, is one.
\end{proof}

\begin{lemma}[Point law and birth intensities]\label{lem:pointlaw}
The formula \(\dd\Prob^{(p)}=Z_p(T)\dd\Prob\) defines a probability
measure. At any bounded stopping time \(\theta\le T\),
\(\E^{(p)}X=\E[Z_p(\theta)X]\) for nonnegative
\(\F_\theta\)-measurable \(X\). Before the root's saving boundary,
the birth intensity of \(q\) under this law is \(J_p(q)\lambda_q\).
\end{lemma}
\begin{proof}
The martingale identity gives
\(\E[Z_p(T)\mid\F_\theta]=Z_p(\theta)\).
One can obtain the stopping-time version from finite-valued dyadic
approximations from above, right continuity and bounded convergence.
For the intensity statement, enumerate the births, condition
\(Z_p(T)\) on each birth's post-jump sigma field, substitute
\(Z_p(t)=J_p(q)Z_p(t-)\), and then apply
\eqref{eq:compensator}. Tonelli justifies the nonnegative birth sum.
On time integrals \(Z_p(t-)=Z_p(t)\) almost everywhere.
\end{proof}

The changed law forbids a birth that kills the root, but does not
delete the self-rate term \(\lambda_p\) from the original deterministic
weight ODE. Losing this distinction changes the compensation identity.

\begin{proposition}[Exact remaining-exposure loss]\label{prop:loss}
With \(E_p\) and \(\theta_p\) as in Section~\ref{sec:overview},
\eqref{eq:previewloss} holds for every fixed \(n\). Therefore
\[
 \sup_p\Prob^{(p)}(E_p)\le\varepsilon
 \quad\Longrightarrow\quad
 \E|S_n|\ge(1-\varepsilon)n A_n.
\]
\end{proposition}
\begin{proof}
Apply \eqref{eq:compensator} to the birth count and use \(\E Z_p(t)=1\):
\[
 \E|S_n|=\frac1n\sum_p\int_0^T m_t(p)\E C_p(t)\dd t
 =nA_n-\frac1n\sum_p\int_0^T m_t(p)\E V_p(t)\dd t.
\]
On a path with positive saving,
\(V_p(t)=Z_p(T)\ind_{\theta_p\le t}\). Tonelli changes its integral
to \(Z_p(T)\int_{\theta_p}^T m_t(p)\dd t\).
Deadline saving has no remaining exposure. A path dying without
positive saving has zero terminal credit and may be assigned
\(\theta_p=T\). The density identity proves \eqref{eq:previewloss}.
Finally the remaining exposure is between zero and the deterministic
total \(F_{\typ(p),T}(p/n)\); sum the probability bound. Events need
not be independent or disjoint.
\end{proof}

\begin{lemma}[Product formula and a positive lower bound]\label{lem:product}
While \(p\) is active,
\begin{equation}\label{eq:product}
 w_p(t)=
 \exp\left\{\int_0^t H_p(u)\dd u+
 \int_0^t\lambda_p(u)\left(1-
       \sum_{\substack{\ell\ni p\\\ell\in\Gamma}}b_{N_\ell}\right)\dd u\right\}
 \prod_{\substack{\ell\ni p\\\ell\in\Gamma}}a_{N_\ell}(k_\ell).
\end{equation}
For large enough \(n\), \(w_p\ge2^{-g-1}\) uniformly while active.
The absolute value of the self-correction integral is at most
\(\max(1,g-1)TB/n=o(1)\).
\end{lemma}
\begin{proof}
Between births, \(\mathrm d\log a_N(k)/\mathrm dt=\Lambda_\ell b_N\).
At a first designated-line birth \(a_1/a_0=1/2\), exactly matching
the surviving point's halving. Subtract these factors from
\eqref{eq:weightflow}; all designated-line terms cancel except the
displayed self-correction. Initially every factor is one. Since
\(H_p\ge0\), \(a_0\ge1\), \(a_1\ge1/2\), and \(\lambda_p\le B/n\),
the asserted bounds follow.
\end{proof}


\section{Mixed moments without independence}\label{sec:moments}

The estimates must retain leading coefficient one. Replacing all
joint probabilities by an unspecified constant multiple of products
would lose the narrow numerical margin. The following compensating
factors separate ordinary-line dependence, designated-line dependence,
and simultaneous changes caused from outside a pair's baseline.

Fix disjoint sets of labels \(C,D\), called roots and witnesses, and
keep their union \(L=C\cup D\) fixed through the argument. If a witness
\(q\) has been born, let \(b_q\) be a deterministic factor
\(e_q(t_q)\in[0,1]\) evaluated at its birth time; otherwise set it to
zero. For a line put \(m_\ell=|C\cap\ell|\). Define
\begin{align}
 \Phi_{C,D}={}&\ind_{C\text{ active}}
 \prod_{q\in D}b_q\prod_{p\in C}w_p\,
 \prod_{\ell\notin\Gamma}r_\ell^{(m_\ell-1)_+}
 \prod_{\ell\in\Gamma}a_{N_\ell}(k_\ell)^{-(m_\ell-1)_+},
 \label{eq:phi}\\
 E_L={}&\prod_{\{a,b\}\subseteq L}\eta_{ab},
 \qquad Y_{C,D}=E_L\Phi_{C,D}.\label{eq:etaall}
\end{align}
Only lines containing two roots contribute nontrivial factors, so the
products are finite. The correction \(E_L\) uses the entire label set,
not just the current roots. This makes it identical on the two sides
of a witness-birth transition.

\begin{lemma}[Generator and source bound]\label{lem:source}
The positive source in the compensated drift of \(Y_{C,D}\) is at
most
\begin{equation}\label{eq:source}
 \sum_{q\in D}\beta_q(t)Y_{C\cup\{q\},D\setminus\{q\}},
\end{equation}
where \(\beta_q=m_t(q)e_q(t)/n\) if no pair of labels in \(L\) shares
a designated line. In general one may multiply each \(\beta_q\) by
\(\max(1,A_1)^g\), where \(A_1=1/(2-K)\).
Every other compensated drift and every auxiliary positive-retirement
jump is nonpositive.
\end{lemma}
\begin{proof}
Fix a birth cause outside the current roots and group roots by the
lines they share with it. On an active ordinary line, the \(m\) root
compensations contribute \(m\lambda_q\), while the factor
\(r^{m-1}\) subtracts \((m-1)\lambda_q\); the group is killed at the
birth. On a designated line of occupancy zero, the continuous
contribution is \(\lambda_q/2\), and the group multiplier is
\((1/2)^m(a_0/a_1)^{m-1}=1/2\). Occupancy one gives continuous
contribution \(\lambda_q\) and multiplier zero. An empty ordinary
line contributes nothing and has multiplier one.

Thus each group has a multiplier \(J_i\in\{0,1/2,1\}\), and its
combined relative drift is
\[
 \lambda_q\left[\sum_i(1-J_i)+\prod_iJ_i-1\right].
\]
For \(d_i=1-J_i\in[0,1]\), the finite Bonferroni inequality gives
\[
 0\le\sum_i d_i-\left(1-\prod_i(1-d_i)\right)
 \le\sum_{i<j}d_i d_j.
\]
One way to prove the upper bound is to introduce independent Bernoulli
variables with success probabilities \(d_i\) just for this finite
algebraic identity: if their sum is \(N\), then
\(N-\ind_{N\ge1}\le\binom N2\), and take expectations.
This auxiliary proof does not assume independence in the geometric
process. Different groups supply a pair of roots whose baseline does
not contain \(q\). The negative logarithmic drift of \(E_L\), namely
\(-\sum_{\{a,b\}\subseteq L}\kappa_{ab}\), pays this upper bound.
There is no jump cross term because \(E_L\) is continuous.
If \(q\) is a root, its own continuous compensation and zeroing jump
cancel; the shared-line correction cancels the other terms on its
line. Thus internal causes introduce no extra positive drift.

It remains to consider an unborn witness \(q\in D\). On a branch
where all roots survive its birth, any shared ordinary line was
previously empty, so its credit factor is one. A shared designated
line containing \(k\ge1\) roots must previously have occupancy zero.
Comparing the post-birth observation with the pre-birth source gives
the exact factor
\[
 2^{-k}a_1^{-(k-1)}a_0^k=a_1.
\]
The rate divided by the source's \(w_q\) is \(m_t(q)/n\).
The witness factor supplies \(e_q(t)\), and \(E_L\) cancels exactly.
There are at most \(g\) shared designated lines. This proves
\eqref{eq:source}. Prior retirement of an auxiliary root only drops
the observation to zero.
\end{proof}

\begin{proposition}[Mixed moments at a common stopping time]\label{prop:mixed}
Write \(R=|L|\) and \(F_q^\beta(t)=\int_0^t\beta_q(u)\dd u\).
On the stopped process,
\begin{equation}\label{eq:mixed}
 \E\Phi_{C,D}(t)\le
 (1-\epsilon_n)^{-\binom R2}\prod_{q\in D}F_q^\beta(t).
\end{equation}
More precisely, for a deterministic \(h\le T\) and any common
bounded stopping time \(\theta\le h\),
\begin{equation}\label{eq:mixedstop}
 \E\Phi_{C,D}(\theta)\le
 (1-\epsilon_n)^{-\binom R2}\prod_{q\in D}F_q^\beta(h).
\end{equation}
The entire fixed-label hierarchy is frozen at the same \(\theta\),
using the postgeometry, prepositive observation convention.
\end{proposition}
\begin{proof}
For \(A\subseteq D\), abbreviate
\[
 Y_A(t)=Y_{C\cup(D\setminus A),A}(t),\qquad
 h_q(t)=\int_t^h\beta_q(u)\dd u.
\]
Every \(Y_A\) uses the same label set \(L=C\cup D\). Define the
nonnegative backward polynomial
\begin{equation}\label{eq:backwardpoly}
 \mathcal P(t)=\sum_{A\subseteq D}Y_A(t)
                  \prod_{q\in D\setminus A}h_q(t),
 \qquad 0\le t\le h.
\end{equation}
For \(q\in A\), the source of \(Y_A\) in Lemma~\ref{lem:source}
is at most \(\beta_qY_{A\setminus\{q\}}\). In the summand indexed
by \(A\setminus\{q\}\), differentiating its factor \(h_q\)
gives exactly the opposite term, since \(h'_q=-\beta_q\).
Thus all positive sources cancel term by term. The remaining
drift and auxiliary-retirement jumps are nonpositive, so
\(\mathcal P\) is a supermartingale. For fixed \(n\) and fixed
finite labels all these stopped observations are bounded, as are
the deterministic coefficients; no unlocalized local-martingale
expectation is being used.

At time zero all summands containing a witness vanish. The remaining
root-only observation is at most one. Hence
\[
 \mathcal P(0)\le\prod_{q\in D}F_q^\beta(h),\qquad
 \E Y_{C,D}(\theta)\le\E\mathcal P(\theta)\le\mathcal P(0).
\]
The first inequality follows from the \(A=D\) term of the
nonnegative sum. Finally divide by the pathwise lower bound
\(E_L\ge(1-\epsilon_n)^{\binom R2}\). This proves
\eqref{eq:mixedstop}; set \(h=\theta=t\) for \eqref{eq:mixed}.
The integrals on the right always end at the deterministic \(h\),
not at the random \(\theta\).

For clarity, saving a boundary observation before simultaneous
positive retirements is implemented by freezing the whole auxiliary
hierarchy at \(\theta\). It does not restore an auxiliary root
that retired earlier. Lemma~\ref{lem:source} already holds after
the geometric birth update without using the subsequent positive
retirement, so it also applies to a birth that triggers this freeze.
At a continuous boundary the values are unchanged by freezing.
\end{proof}

For fixed \(R\), the first factor in \eqref{eq:mixed} is \(1+o(1)\).
Removing connection factors costs only a constant depending on
\(R,g,K,\rho\), because \(r\ge\rho\) and \(a_N\le2/(2-K)\).
In particular, products with \(j\) distinct born witnesses cost
at most \(C_R(t/n)^j\). A repeated root weight is bounded by
\(w^k\le B^{k-1}w\). A witness repeated many times is still just one
birth and must not be charged as independent births.

For estimates under \(\Prob^{(p)}\), put \(p\) among the roots of
every member of the hierarchy, and freeze all of it at the first
positive root boundary or global stop, capped at \(T\). Use the
postgeometry, prepositive observation convention. Lemma~\ref{lem:pointlaw}
then converts the expectation to the desired point-law moment.
No witnesses born after this common boundary may be added.
These conventions are needed for estimates at first crossings,
not just at deterministic times. In the witness-free case the
polynomial is the single observation itself. Its supermartingale
property therefore also supplies conditional estimates between
two observation times, as used below to remove ordinary-line credit.


\section{Removing the global safety stops}\label{sec:global}

All estimates here are proved for the stopped process. On it,
\(k\le K\), \(r\ge\rho\), and \(\eta\ge1-\epsilon_n\), so
Proposition~\ref{prop:mixed} is available before any probability bound
on the global stop has been established.

\subsection{Pair-correction and virtual-deletion guards}
\begin{lemma}[Pair-correction tail]\label{lem:etatail}
For each fixed positive integer \(m\),
\begin{equation}\label{eq:etatail}
 \Prob(\text{pair-correction guard})
 \le C_m n^{4-m+2m\xi+d(m-1)+m\zeta}.
\end{equation}
\end{lemma}
\begin{proof}
For a fixed pair \(a,b\), classify a cause changing both roots
according to whether the two lines are ordinary or designated.
Two ordinary lines require two distinct selected witnesses; one
ordinary line requires one; two designated lines require only an
intersection of two among finitely many directions.

In an \(m\)-fold expansion, regard distinct selected witnesses as
vertices and the double-ordinary causes as edges of a multigraph.
If there are \(j\) witnesses and \(c\) connected components, the
first edge of a component costs \(C_\xi n^{2+\xi}\) by
\eqref{eq:twostar}. An edge adding one new witness costs
\(C_\xi n^{1+\xi}\) by \eqref{eq:fixedline}. With both witnesses
known, the cause is the intersection of two fixed lines and has at
most one choice. Coincidence of those lines would put the cause on
the excluded baseline. Thus the count is at most
\(C_m n^{j+(j-c)\xi}\le C_m n^{j+2m\xi}\).
A component containing a designated/ordinary cause may be started
with that cause, at cost \(C_\xi n^{1+\xi}\). A double-designated
cause has at most \(g^2\) choices. This also covers mixed components
and a witness appearing on different sides of different edges.

If the physical causes have \(k\ge1\) distinct labels, their rates
supply \(n^{-m}\), root repetitions cost
\(B^{m-k}\le B^{m-1}\), and the \(j\) distinct witnesses cost
\(C_m n^{-j}\) by Proposition~\ref{prop:mixed}. Root activity
indicators can be replaced by constant multiples of root weights
using Lemma~\ref{lem:product}. A label that is both an active cause
and a selected witness gives zero. Therefore
\[
 \E\kappa_{ab}(t)^m\le C_m n^{-m+2m\xi}B^{m-1}.
\]
Hitting the floor requires
\(\int_0^T\kappa_{ab}\dd t\ge-\log(1-\epsilon_n)\ge\epsilon_n\).
Apply Holder to the time integral, then Markov and a union bound
over at most \(n^4\) pairs. This is \eqref{eq:etatail}.
\end{proof}

\begin{lemma}[Virtual-deletion tail]\label{lem:virtualtail}
For each fixed positive integer \(m\),
\begin{equation}\label{eq:virtualtail}
 \Prob(\text{virtual-deletion guard})
 \le C_m\sum_{a=1}^m n^{3-\gamma m+d(m-a)+a\xi}.
\end{equation}
\end{lemma}
\begin{proof}
Fix a designated \(\ell\) and an active \(q\notin\ell\).
In an \(m\)-fold expansion of \eqref{eq:virtual}, suppose there are
\(a\) distinct leaves on \(\ell\). The lines from \(q\) to these
leaves are distinct, and their selected witnesses are distinct:
a common witness on two such lines would be \(q\), which is active.
By \eqref{eq:fixedline} the number of leaf/witness choices is at most
\(C_m n^{a(1+\xi)}\). Source prices supply \(n^{-a}\), and repeated
leaves cost \(B^{m-a}\). The normalized weight is at most \(2w\).
Consequently
\[
 \E[\ind_{q\text{ active}}\mathcal V_{\ell,q}^m]
 \le C_m\sum_{a=1}^m B^{m-a}n^{a\xi}.
\]
Apply the same fixed-label hierarchy at the first crossing; no
union over observation times is necessary. Specifically, include
\(q\) and all distinct leaves as roots, the selected labels as
witnesses, and apply \eqref{eq:mixedstop} with deterministic
terminal time \(T\). The activity indicator of \(q\) is bounded
by a constant times its positive weight by Lemma~\ref{lem:product}.
There are \(O(n)\)
designated lines and \(n^2\) possible external points. Markov at
level \(n^\gamma\) gives \eqref{eq:virtualtail}.
\end{proof}

The slope in \(m\) in \eqref{eq:etatail} is
\(-(1-d-2\xi-\zeta)<0\). Since \(\xi<d\), the largest exponent in
\eqref{eq:virtualtail} occurs at \(a=1\) and equals
\(3-d+\xi-m(\gamma-d)\). Both probabilities are therefore smaller
than any prescribed fixed power of \(n^{-1}\), after choosing one
sufficiently large fixed moment order. The process itself is unchanged
by that choice of proof order.

For the concrete parameters
\(d=19999/20000\), \(\gamma=39999/40000\), and
\(\xi=\zeta=10^{-6}\), one may already take
\(m=2^{18}=262144\). Including the factor \(B\) for a change to
any point law, the two exponents are respectively
\[
 4-m+2m\xi+d(m-1)+m\zeta+d=-8.320768,
 \qquad 3-\gamma m+d(m-1)+\xi+d=-3.553599.
\]
Only their negativity is needed. The constants depending on this
fixed order need not be evaluated to obtain an asymptotic bound.

\subsection{The protected law of a designated line}
Fix a designated line and one of its occupancy stages \(N=0,1\).
Starting at the stage's entry time \(U\), suppress its next birth
through the density process
\[
 D_t=\ind_{\text{no new line birth in }(U,t]}
 \exp\left(\int_U^t\Lambda_\ell(s)\dd s\right),
\]
stopped and saved at the relevant guards or \(T\). Before \(U\) its
value is one. The birth at \(U\), if any, starts the stage and is
not counted again as its next birth. The continuous drift
\(D\Lambda_\ell\) cancels the zeroing-jump drift. Also
\[
 \int_U^t\Lambda_\ell\dd s
 =\int_{k(U)}^{k(t)}a_N(k)\dd k
 \le2\log\frac2{2-K}.
\]
Thus \(D\) is a bounded martingale and defines a probability law
conditionally on the entry history. Its upper bound is the fixed
constant \(A_0^2\), where \(A_0=2/(2-K)\). On a crossing occurring
before the next line birth, \(D\ge1\), so the original crossing
probability is at most its probability under the protected law.
In particular this comparison does not introduce the upper-density
factor \(A_0^2\). Under the protected law the next line birth has
rate zero; the external birth rates remain their original rates.

Normalize each active line weight by
\(c_p=w_p/a_N(k_\ell)\), and save this \emph{normalized} value at
positive retirement. It is not subsequently divided by a changing
\(a_N\). For a fixed cell/type subset of the line, let \(M\) be the
sum of current and saved \(c_p\). Define its entropy sum using
\[
 h(c)=c\left(1+\log\frac{2B}{c}\right),\quad h(0)=0,
 \qquad H=\sum h(c_p),
\]
again with saved terms frozen. Under the protected law the surviving
positive drift of each credit is
\(g_p=\lambda_p(1-b_N)c_p\). This follows by subtracting
\(\mathrm d\log a_N/\mathrm dt=\Lambda_\ell b_N\) from the original
ODE and deleting only the line-birth jump rates. It is
\(\lambda_pc_p/2\) at occupancy zero and zero at occupancy one.

\begin{lemma}[Credit, entropy and whole-event energy]\label{lem:entropy}
Stop further at \(M\le C_0n\). There are compensated martingales
\(Y_M,Y_H\) and a nonnegative cumulative cost \(V\) such that
\begin{align}
 M&=M(U)+D_M+Y_M,&D_M&=O(B),\label{eq:creditdecomp}\\
 H&=H(U)+D_H-V+Y_H,&D_H&=O(B\log n),\label{eq:entropydecomp}
\end{align}
and, before the virtual guard,
\begin{align}
 \langle Y_M\rangle&\le C n^\gamma V,
 &|\Delta Y_M|&\le C n^\gamma,\label{eq:bracketM}\\
 \langle Y_H\rangle&\le C n^\gamma(\log n)^2V,
 &|\Delta Y_H|&\le C n^\gamma\log n.\label{eq:bracketH}
\end{align}
The constants are independent of the admissible entry history.
\end{lemma}
\begin{proof}
For a multiplier \(J\in\{0,1/2,1\}\), direct algebra gives
\[
 h'(c)c(1-J)+h(Jc)-h(c)=-c(J\log J-J+1).
\]
The cost is \(c\) for deletion, \(c(1-\log2)/2\) for halving,
and zero otherwise. It is comparable by fixed positive constants
to the credit decrease \(c(1-J)\). Define \(V\) to be the rate
integral of the \emph{sum of these costs for each complete birth
event}. Compensating gives \eqref{eq:creditdecomp} and
\eqref{eq:entropydecomp}. Since \(\lambda_p\le B/n\) and the
credit sum is capped by \(C_0n\), the first drift is \(O(B)\).
The active credits have the positive lower bound supplied by
Lemma~\ref{lem:product} divided by \(A_0\), so
\(h'(c)=\log(2B/c)=O(\log n)\) on them. This gives the entropy
drift estimate.

In one external birth event, the ordinary-line deletion of normalized
mass is at most \(M_n\), and effects through other designated
directions touch at most \(g\) points, each of credit at most \(2B\).
Its total credit decrease \(\Delta\) is therefore at most
\(M_n+2gB=O(n^\gamma)\). The square of the \emph{whole} decrease
satisfies \(\Delta^2\le C n^\gamma\Delta\).
The absolute entropy decrease is at most \(C\log n\) times this
credit decrease, including deletion to zero. Comparing the decrease
to the event's total cost and integrating its birth rate proves
\eqref{eq:bracketM}--\eqref{eq:bracketH}. These estimates retain
all cross terms of simultaneous changes. Summing the squares of
individual point changes would not prove the required bracket bound.
The birth that causes a global stop is included: its decrease is
controlled by the virtual-guard bound in the state immediately
before that birth. Positive saving itself changes neither the
credit ledger nor its entropy.
\end{proof}

\begin{lemma}[Exponential crossing estimate]\label{lem:exp}
Let \(Y\) be a compensated birth martingale with \(Y_0=0\), jumps
bounded in absolute value by \(b\), and continuous predictable
quadratic variation \(Q\). For \(x,v>0\),
\begin{equation}\label{eq:freedman}
 \Prob(\exists t\le T:Y_t\ge x,\ Q_t\le v)
 \le\exp\left(-\frac{x^2}{2(v+bx/3)}\right).
\end{equation}
\end{lemma}
\begin{proof}
For \(0<\theta b<3\) and \(|u|\le b\), the Taylor series and
\(k!\ge2\cdot3^{k-2}\) for \(k\ge2\) give
\[
 e^{\theta u}-1-\theta u
 \le\psi u^2,\qquad \psi=\frac{\theta^2}{2(1-\theta b/3)}.
\]
Stop at the first \(Y\ge x\), capped at \(T\). The stopped
exponential \(\exp(\theta Y-\psi Q)\) is bounded by
\(\exp(\theta(x+b))\). Its compensated drift is nonpositive by the
displayed inequality, so its expectation is at most one.
Since \(Q\) is continuous and nondecreasing, a crossing with
\(Q\le v\) exists precisely when the first crossing has \(Q\le v\).
On this event the terminal stopped exponential is at least
\(\exp(\theta x-\psi v)\). Markov gives
\(\exp(-\theta x+\psi v)\); take \(\theta=x/(v+bx/3)\).
This proof uses a bounded stopped exponential, rather than assuming
that an unlocalized exponential local martingale has mean one.
\end{proof}

\begin{proposition}[Clock stops are rare]\label{prop:clock}
For each fixed positive tolerance and admissible fixed data, the
failure probability of the simultaneous designated-line credit
bound is at most
\[
 Cn\exp\left(-c\frac{n^{1-\gamma}}{\log n}\right).
\]
In particular, the probability of a designated clock reaching \(K\)
is smaller than every fixed inverse power of \(n\).
\end{proposition}
\begin{proof}
Stop the cumulative cost at \(C_1n\log n\). Since \(H\ge0\),
\(H(U)=O(n\log n)\), and \(D_H=o(n\log n)\), reaching this cost
cap, for sufficiently large \(C_1\), forces an upward fluctuation
of \(Y_H\) of order \(n\log n\). Apply Lemma~\ref{lem:exp} with
\(v=O(n^{1+\gamma}(\log n)^3)\) and
\(b=O(n^\gamma\log n)\). An increase of \(M\) by a fixed multiple
of \(n\) similarly forces a fluctuation of \(Y_M\) of order \(n\),
with \(v=O(n^{1+\gamma}\log n)\) and \(b=O(n^\gamma)\).
Both exponents have order \(n^{1-\gamma}/\log n\).

At the first birth on the line,
\((w/2)/a_1=w/a_0\); the selected point is removed. The saved ledger
is unchanged and the moving normalized credit cannot increase.
At the second birth, all remaining moving credit disappears. To
treat the occupancy-one stage, condition on its actual entry history
and the preceding good credit bound. At time zero the ledger of a
cell/type piece equals its lattice-point count \(L_E\le n\).
On a good occupancy-zero history it is at most
\(L_E+\varepsilon n\), and the first line birth cannot increase it.
The lower and upper credit bounds from Lemma~\ref{lem:product}
give \(h(c)\le C(\log n)c\), including saved positive credits.
Consequently both initial credit and entropy bounds hold uniformly
at the occupancy-one entry time. Apply the same conditional bound
to every such history and average it. The union is only over the
two stages, not over possible histories or starting times.

There are \(O(n)\) designated lines and a fixed number of cell/type
pieces. The estimates control whole time intervals, so no separate
union over phases is needed. Summing their crossing probabilities preserves
the displayed bound. On the resulting good event, integrate the
moving normalized credit with the masks to obtain
\[
 k_\ell\le\text{deterministic line capacity}
             +C_{\mathrm{data}}T\varepsilon+O(n^{-1}).
\]
Choose \(\varepsilon>0\) first so that the second term is smaller
than half of \(\kappa_H\), then choose \(n\) large. The capacity
condition \eqref{eq:heavyinput} precludes hitting \(K\).
Positive retirement need not be rare for this argument: its credit
was retained in the ledger throughout.
\end{proof}

\subsection{Explicit clock constants for the verified input}

For the input used in the application,
\[
 g=16,\quad K=255/128,\quad A_0=256,\quad G=256,
 \quad T=\frac{30088384743729}{25000000000}<2^{11}.
\]
Write \(C_{\rm line}=4G=1024\). The verified designated-line
capacity bound has margin
\[
 \kappa_H=\frac{69405753}{80000000000},\qquad
 \varepsilon_H=\frac{\kappa_H}{8TC_{\rm line}}
       =\frac{115676255}{1314581588376682496}>2^{-34}.
\]
The following explicit estimate is not needed to compute a final
threshold for all five stopping causes. It records how the
history-uniform clock argument applies to this particular input.

\begin{corollary}[A quantitative clock budget]\label{cor:clockexplicit}
For all sufficiently large \(n\),
\begin{equation}\label{eq:clockexplicit}
 \sup_p\Prob^{(p)}(\text{clock is an early cause})
 \le221184\,n^{1+d}
 \exp\left(-\frac{\varepsilon_H^2 n^{1-\gamma}}
                     {2^{18}\log n}\right).
\end{equation}
In fact, with \(d=19999/20000\) and \(\gamma=39999/40000\),
\begin{equation}\label{eq:clockthreshold}
 \log n\ge10^7\quad\Longrightarrow\quad
 \sup_p\Prob^{(p)}(\text{clock is an early cause})<2\cdot10^{-6}.
\end{equation}
\end{corollary}
\begin{proof}
Assume first \(n^{1-d}\ge30T\). The self-correction in
\eqref{eq:product} has absolute value at most \(15TB/n\le1/2\),
so \(w_p\ge2^{-17}\) while active. Every positive normalized
credit, including a saved credit, therefore satisfies
\[
 2^{-25}\le c_p\le2B.
\]
For \(L_n=32\log n\) and \(n\ge2\), this implies
\(h(c)\le L_nc\) and \(0\le h'(c)\le L_n\) on positive active
credits. Stop temporarily while \(M\le3n\). The drift terms in
Lemma~\ref{lem:entropy} obey
\[
 0\le D_M\le\tfrac32TB,\qquad 0\le D_H\le L_nD_M.
\]
For a complete external birth let \(v\) be its entropy cost and
\(\Delta\) its credit loss. Since \(\log2<3/4\), deletion and
halving both satisfy \(\Delta/4\le v\le\Delta\). The virtual
guard and at most 16 other designated-line intersections give
\(\Delta\le n^\gamma+32B\le33n^\gamma\). Hence
\begin{align*}
 \langle Y_M\rangle&\le132n^\gamma V,
 &|\Delta Y_M|&\le33n^\gamma,\\
 \langle Y_H\rangle&\le132n^\gamma L_n^2V,
 &|\Delta Y_H|&\le33n^\gamma L_n.
\end{align*}
The bound on an entropy jump follows directly from
\(h(c)-h(Jc)\le L_nc(1-J)\), including \(J=0\);
it does not assume a bounded derivative at zero.

Consider a stage starting with \(M(U)\le2n\), and any
\(0<\varepsilon\le1\) for which
\(n^{1-d}\ge3T/\varepsilon\). Stop at the first of a credit
increase \(\varepsilon n\), cost \(V=4L_nn\), an original guard,
and \(T\). Credit increases and \(V\) are continuous, so before
this stop \(M\le3n\), and \(D_M\le\varepsilon n/2\).
If the cost boundary is reached first, \(H\ge0\),
\(H(U)\le2L_nn\) and \(D_H\le L_n\varepsilon n/2\) force
\(Y_H\ge L_nn\). Lemma~\ref{lem:exp}, with
\(v=528n^{1+\gamma}L_n^3\) and \(b=33n^\gamma L_n\), bounds
its probability by
\(\exp[-n^{1-\gamma}/(1078L_n)]\).
If the credit boundary is reached first, \(Y_M\ge\varepsilon n/2\)
and \(\langle Y_M\rangle\le528L_nn^{1+\gamma}\), giving
\(\exp[-\varepsilon^2n^{1-\gamma}/(4268L_n)]\).
As \(4268\cdot32<2^{18}\), their sum is at most
\[
 2\exp\left(-\frac{\varepsilon^2n^{1-\gamma}}
                         {2^{18}\log n}\right).
\]
This bound is conditional on the entry history and its constants
are independent of that history.

At the first line birth the ledger cannot increase. Applying the
last estimate first at occupancy zero and then conditionally at
occupancy one shows that
\(M_E(t)\le L_E+2\varepsilon n\) fails with at most twice the
displayed probability. The 16 primitive unoriented directions of
height at most three satisfy
\(\sum_{v\in\Gamma}(|v_1|+|v_2|)=54\), so they determine at
most \(54n\) lines. Each line has at most 1024 nonempty cell/type
pieces, by Lemma~\ref{lem:cells} and the corresponding axis count.
Thus the simultaneous ledger bound fails under the original law
with probability at most
\[
 221184\,n\exp\left(-\frac{\varepsilon^2n^{1-\gamma}}
                                  {2^{18}\log n}\right).
\]

On its complement, integration against the masks gives
\[
 k_\ell(t)\le\frac1n\sum_{q\in\ell}F_{\typ(q),T}(q/n)
                   +2\varepsilon TC_{\rm line}.
\]
For this input the field bound is 7 and there are \(256^2\)
common rectangles. The lattice-to-chord argument without removing
the root has error at most
\(E_{\rm cap}/n\), where \(E_{\rm cap}=4\cdot256^2\cdot7=1835008\).
Take \(\varepsilon=\varepsilon_H\) and
\(n\ge4E_{\rm cap}/\kappa_H\). The lattice error and ledger
error are each at most \(\kappa_H/4\), so \(k_\ell<K\).
Finally \(Z_p(T)\le B\) changes to any point law at cost at most
\(B=n^d\), proving \eqref{eq:clockexplicit}.

To verify the stated threshold, put \(s=\log n\ge10^7\).
The rational bounds \(\varepsilon_H>2^{-34}\), \(T<2^{11}\),
\(\kappa_H>2^{-11}\), and \(E_{\rm cap}<2^{21}\) imply all
three preliminary conditions, since \(n^{1-d}=e^{s/20000}\ge
e^{500}>2^{500}\). The right side of \eqref{eq:clockexplicit}
is at most
\[
 2^{18}\exp\left(2s-\frac{e^{s/40000}}{2^{86}s}\right).
\]
The function \(e^{s/40000}/s^2\) increases for \(s\ge80000\).
At \(s=10^7\), use \(e^{250}>2^{250}\) and
\(10^{14}<2^{47}\) to obtain
\(e^{s/40000}/(2^{86}s)>4s\); this persists for larger \(s\).
The result is bounded by \(2^{18}e^{-2s}<2\cdot10^{-6}\).
This is only a threshold for the clock cause, not a value of the
eventual threshold in the final existence theorem.
\end{proof}

Combining the three guards shows, for every fixed \(A_0'>0\),
\begin{equation}\label{eq:globaltail}
 \Prob(\sigma\le T)\le C_{A_0'}n^{-A_0'}.
\end{equation}
Under any point law this is at most \(B\) times as large, hence is
\(o(1)\) uniformly in the root. This establishes global-stop smallness
before the weight-cap and ordinary-line tag estimates below.




\section{Ordinary-line risk under the point law}\label{sec:riskproof}

Fix a root \(p\) and its deterministic deadline \(\tau=\tau_p\).
All observations below are stopped and saved at the root's first
positive boundary or global stop. A quantity explicitly restricted
to the active root is zero afterwards. Write
\(F_q(t)=\int_0^t m_u(q)\dd u\), and let
\[
 I_p=\int_0^{\theta_p\wedge T}\ind_{p\text{ active}}H_p(t)\dd t
\]
be its integrated ordinary-line hazard. Earlier root death has
\(\Prob^{(p)}\)-probability zero. We will bound \(I_p\) by a proxy
having sharp fixed-order moments, plus an error small in \(L^1\).
The distinction between those two kinds of estimate is essential.

\subsection{Removing ordinary-line credit at additive cost}
For an ordinary line through \(p\), let
\(H_\ell(t)=r_\ell(t)\sum_{q\in\ell\setminus\{p\},\ q\text{ active}}w_q(t)\).
The two-root, no-witness case of Proposition~\ref{prop:mixed}, followed
by the conditional density identity, gives
\begin{equation}\label{eq:conditionalH}
 \E^{(p)}[\ind_{p\text{ before boundary at }t}H_\ell(t)\mid\F_u]
 \le(1-\epsilon_n)^{-1}
       \ind_{p\text{ before boundary at }u}H_\ell(u)
 \quad(u\le t).
\end{equation}
For each leaf the witness-free observation
\(\eta_{pq}w_pr_\ell w_q\) is a supermartingale. Under the
point law this bounds the conditional expectation of
\(\eta_{pq}r_\ell w_q\); the ratio of its correction factors is
at most \((1-\epsilon_n)^{-1}\). Conditional division is performed only
where \(Z_p(u)>0\); its complement is null under the point law.
The auxiliary-root sum excludes \(p\), so no self-square is
silently included in this two-root argument.

One selected witness \(z\in\ell\) and two leaves give
\begin{equation}\label{eq:secondH}
 \E^{(p)}[\ind_{p\text{ before boundary},\ z\in S_t}H_\ell(t)^2]
 \le C L_\ell(L_\ell+B)t/n.
\end{equation}
For distinct leaves the connection factor is exactly \(r_\ell^2\);
a repeated leaf is paid using \(r^2w^2\le Brw\).
The one-leaf version is at most \(CL_\ell t/n\).
All labels here lie on the same ordinary line, so there is no
designated-line source multiplier.

Since \(r\) starts moving at the first line birth,
\(1-r_\ell(t)=\int_{\text{activation}}^t r_\ell\Lambda_\ell\dd u\).
The correct self-rate bound is
\begin{equation}\label{eq:selfrate}
 r_\ell\Lambda_\ell\le H_\ell/n+B/n.
\end{equation}
Multiplying the leaf sums in the hazard by \(r_\ell\) therefore loses
a nonnegative amount whose expectation is at most
\begin{equation}\label{eq:removecredit}
 Cn^{-3}\sum_{\ell\ni p}(L_\ell^3+BL_\ell^2)
 +CBn^{-3}\sum_{\ell\ni p}L_\ell^2
 =O(1+(B/n)\log n)=o(\log n).
\end{equation}
For completeness, insert the integral for \(1-r\), use
\eqref{eq:conditionalH} to move the future \(H\) to the earlier
time, cover activation by \(\sum_{z\in\ell}\ind_{z\in S_u}\),
and apply \eqref{eq:secondH}. The second term in
\eqref{eq:removecredit} is the separate self-rate contribution from
\eqref{eq:selfrate}. The line sums are Lemma~\ref{lem:lines}.
This is an additive error, not a constant-factor loss in the leading
\(\log n\) coefficient.

\subsection{Long lines: replace the leaf sum, retain the birth time}
Call a line long when \(L_\ell\ge B\log n\). Fix a selected witness
\(z\in\ell\), and put
\[
 H=r_\ell\sum_{q\in\ell\setminus\{p\},\ q\text{ active}}m_t(q)w_q,
 \qquad D_z(t)=\sum_{q\in\ell\setminus\{p,z\}}m_t(q).
\]
For every fixed integer \(k\), the mixed hierarchy gives
\begin{equation}\label{eq:longmoment}
 \E^{(p)}[\ind_{p\text{ before boundary},\ z\in S_t}H^k]
 \le(1+o(1))\frac{F_z(t)}n
       [D_z(t)^k+C_kB L_\ell^{k-1}].
\end{equation}
Indeed, the all-distinct leaves have exactly the required connection
\(r_\ell^k\) and unit source price. A term with \(a<k\) distinct
leaves costs \(B^{k-a}L_\ell^a\le BL_\ell^{k-1}\), since
\(B\le L_\ell\). There are finitely many repetition patterns for
fixed \(k\). No division by a positive lower bound on \(F_z\) is
used; if it is zero the corresponding birth event is null.

For fixed \(\epsilon>0\), use
\(H\ind_{H>D+\epsilon L}\le H^k/(D+\epsilon L)^{k-1}\).
Since \(D\le L\), the main term is at most
\(L(1+\epsilon)^{-(k-1)}\), while the repeated-leaf term is
\(C_{k,\epsilon}B\). Sum over witnesses, long lines and time,
with the hazard factor \(1/n\). As \(B/L\le1/\log n\), the
positive error \(R_{\rm long}\) made by replacing the leaf sum by
\(D_z(t)\) satisfies
\begin{equation}\label{eq:longerror}
 \limsup_{n\to\infty}\sup_p
 \frac{\E^{(p)}R_{\rm long}}{\log n}
 \le C\epsilon+C(1+\epsilon)^{-(k-1)}.
\end{equation}
First choose \(\epsilon\), next a fixed sufficiently large \(k\),
and only then let \(n\) grow. This proves an \(o(\log n)\)
\(L^1\) error for the same process.

For a witness born at time \(u\), define its deterministic remaining
cost and birth price by
\[
 A_z(u)=\frac1n\int_u^\tau D_z(t)\dd t\le\frac{TL_{pz}}n,
 \qquad b_z=\frac1n\int_0^\tau m_u(z)A_z(u)\dd u.
\]
Let \(P_{\rm long}\) be the sum of \(A_z\) over long-line witnesses
born before the root boundary, and let \(M_{\rm long}=\sum_z b_z\).
Insert \(A_z(u)/(TL_{pz}/n)\le1\) as a deterministic birth-retention
factor in the mixed hierarchy, and use \eqref{eq:mixedstop}
with deterministic horizon \(h=\tau\). The common stopping time
limits which witnesses are included; it does not replace \(\tau\)
inside the deterministic cost or birth-price integral. If \(T=0\)
there are no witnesses and the assertion is immediate.
It follows that for every fixed \(m\),
\begin{equation}\label{eq:longproxy}
 \|P_{\rm long}\|_{L^m(\Prob^{(p)})}
 \le M_{\rm long}+o(\log n).
\end{equation}
Here are the non-generic terms in this assertion. A repeated witness
block of size \(r\ge2\) has total price at most
\(C_r n^{-(r+1)}\sum_\ell L_\ell^{r+1}=O_r(1)\), whereas a
singleton block costs \(O(\log n)\). Thus repetitions lose a power
of \(\log n\). For two distinct witnesses sharing a designated
line, that line cannot pass through the root. By
\eqref{eq:fixedline}, the weighted configuration count is
\[
 \sum_{z_1}L_{pz_1}
 \sum_{\substack{z_2\text{ on a designated}\\\text{line through }z_1}}L_{pz_2}
 \le C n^{3+\xi}\log n.
\]
The two birth-and-cost factors supply \(n^{-4}\), giving
\(O(n^{-1+\xi}\log n)\). Remaining singleton factors cost
\(O(\log n)\) each. A finite union over pairs covers multiple
exceptions. The all-distinct, unshared case keeps the exact product
of the \(b_z\), with leading factor \(1+o(1)\).

\subsection{Short stars and their exceptional configurations}
An active short line through \(p\) has one selected witness \(z\).
Each active leaf \(q\ne p\) on that line gives a star \((p,q,z)\).
Let
\[
 X_t=\sum_{\text{short stars}}m_t(q)r_{\ell(p,q)}w_q,
 \qquad
 D_t=\frac1n\sum_{\text{deterministic short stars}}m_t(q)F_z(t).
\]
The actual sum is zero after the root boundary. In the deterministic
sum, \(p,q,z\) are distinct and collinear on a short ordinary line;
there is no selected-point condition.
Put \(t_0=n^{\gamma-1}\).

\begin{lemma}[Short-star moment estimate]\label{lem:short}
For each fixed positive integer \(m\), uniformly for \(t_0\le t\le T\)
and all roots,
\begin{equation}\label{eq:shortmoment}
 \E^{(p)}X_t^m\le(1+o(1))D_t^m+\epsilon_m(n)(tn\log n)^m,
 \qquad \epsilon_m(n)\longrightarrow0.
\end{equation}
More explicitly, for all sufficiently large \(n\) (depending on
the fixed \(m\)), one may use
\begin{equation}\label{eq:shortexplicit}
 \E^{(p)}X_t^m\le[1+C_m n^{-\zeta}]D_t^m+
 C_m\bigl(n^{d-\gamma}+n^{\gamma-1}+n^{-1+\xi}\bigr)
                (tn\log n)^m.
\end{equation}
\end{lemma}
\begin{proof}
In an \(m\)-fold product let \(a\) be the number of distinct leaves
and \(j\) the number of distinct witnesses. The active root makes
the selected witness on each of its lines unique, so
\(1\le j\le a\le m\). A leaf coinciding with a witness or the root
gives zero. If \(j<m\), the first leaf and witness of each star
line have \(O(n^2\log n)\) choices by \eqref{eq:line2}.
Additional leaves on that line have at most \(B\log n\) choices,
and repeated weights cost \(B^{m-a}\). The source price is
\(C_m(t/n)^j\). Dividing this bound by \((tn\log n)^m\) gives
\[
 C_m\left(\frac B{nt}\right)^{m-j}(\log n)^{a-m}=o(1),
\]
uniformly since \(B/(nt)\le n^{d-\gamma}\).

It remains to treat \(j=a=m\), with distinct star lines. The
following list exhausts the extra connections among different stars.
\begin{enumerate}[label=(\roman*)]
\item Two labels share a designated line. Fix the first label.
The line is outside the root, so \eqref{eq:fixedline} bounds the
second label together with its star partner by \(C n^{1+\xi}\),
instead of the generic \(O(n^2\log n)\) row. This covers leaf--leaf,
leaf--witness and witness--witness sharing.
\item A line connecting two leaves has length greater than
\(n^\gamma\). For a fixed first leaf there are
\(O(n^{2-\gamma})\) possible second leaves: sum the height counts
up to \(O(n^{1-\gamma})\). Its star witness has at most
\(B\log n\) choices. Relative to a generic star row the saving
is \(O(n^{d-\gamma})\).
\item A leaf--leaf line of length at most \(n^\gamma\) is active
through a new witness. There are at most \(n^\gamma\) positions
for it and an additional source price \(C_mt/n\), giving
\(O(n^{\gamma-1})\).
\item That line is active through an existing third star witness.
It cannot be either leaf's own witness, since then their star lines
would coincide. For a third witness use
\(\sum_{z\in\ell(q_i,q_j)}L_{pz}\le C n^{1+\xi}\), keeping all
\(m\) original birth prices. This again saves \(n^{-1+\xi}\).
An active leaf or the root cannot itself be a selected witness.
\end{enumerate}
Outside these cases the additional ordinary-line credits are all
one, there are no designated-line connections, and the star credits
exactly match the mixed-moment factors. The source price is the
exact product of \(F_z(t)/n\); keeping the masks produces \(D_t^m\).
Each single exception is bounded after dropping other restrictions,
so a finite union over label pairs also covers overlapping exceptions.
Removing remaining factors costs only \(C_m\), including
\(\rho^{-O(m^2)}\). The savings above and the correction
\((1-\epsilon_n)^{-O(m^2)}=1+O_m(n^{-\zeta})\) prove
\eqref{eq:shortexplicit} and hence the assertion.
\end{proof}

For the parameters used here, \(d-\gamma=\gamma-1=-1/40000\)
and \(-1+\xi=-999999/1000000\), while \(\zeta=10^{-6}>0\).
All these errors vanish for every fixed order. No uniform control
of \(C_m\) as \(m\) grows is asserted or used.

Minkowski's inequality now yields
\begin{equation}\label{eq:shortintegral}
 \left\|\int_{t_0}^\tau X_t/n\dd t\right\|_m
 \le\int_{t_0}^\tau D_t/n\dd t+o(\log n),
\end{equation}
with the integral zero when \(\tau<t_0\).
For earlier times the one-witness, one-leaf bound gives
\(\E^{(p)}X_t=O(tn\log n)\). Its integrated contribution is
\(O(t_0^2\log n)=o(\log n)\) in \(L^1\). No higher moment for
this early-time error is asserted.

\subsection{Hazard concentration and the weight cap}
Collect the credit-removal error, long-line leaf replacement error,
and early short-star contribution into \(R_n\ge0\), so
\(\sup_p\E^{(p)}R_n=o(\log n)\). Define
\[
 Y_n=P_{\rm long}+\int_{t_0}^\tau X_t/n\dd t.
\]
Then \(I_p\le Y_n+R_n\), and \eqref{eq:longproxy} and
\eqref{eq:shortintegral} give, for every fixed \(m\),
\[
 \|Y_n\|_m\le M_{\rm long}+M_{\rm short}+o(\log n)
 \le\mathcal B_n^{\Gamma^c}(p;F_\tau)+o(\log n).
\]
The last inequality is \eqref{eq:timesym}: the division into long
and short lines is symmetric in leaf and witness, and all four
fields use the same \(\tau\).
Proposition~\ref{prop:pair} and \eqref{eq:riskinput} therefore give
the following summary, uniformly in the root and for every fixed
positive integer \(m\):
\begin{equation}\label{eq:hazarddecomp}
 I_p\le Y_n+R_n,\qquad
 \|Y_n\|_{L^m(\Prob^{(p)})}\le(r_*+o_m(1))\log n,
 \qquad \E^{(p)}R_n=o(\log n).
\end{equation}

For any \(b>r_*\), choose \(\epsilon>0\) with
\(r_*+\epsilon<b-\epsilon\). Markov on \(R_n\), followed by
the fixed-\(m\) moment bound, shows
\[
 \limsup_{n\to\infty}\sup_p\Prob^{(p)}(I_p>b\log n)
 \le\left(\frac{r_*+\epsilon}{b-\epsilon}\right)^m.
\]
The right side tends to zero as the fixed proof order \(m\) is
chosen larger. This quantifier order does not require constants
uniform over all \(m\), or an order depending on \(n\).
By \eqref{eq:product}, hitting \(B=n^d\) requires
\(I_p\ge d\log n-O(1)\). Choose \(r_*<b<d\). We have proved
\begin{equation}\label{eq:captail}
 \sup_p\Prob^{(p)}(\text{weight cap is an early saving cause})=o(1).
\end{equation}


\section{Ordinary-line tags and the final loss estimate}\label{sec:closing}

The ordinary-line tag \(r_\ell=\rho\) must be controlled uniformly
over all directions. We use fixed moment orders throughout this
section. Retaining the line's height and counting only the spatial
cells that it meets avoids a growing moment order or constants
uniform over all orders.

\begin{lemma}[Linewise cell count]\label{lem:cells}
A nonaxis line meets at most \(2G-1\) open cells of a fixed
\(G\times G\) grid. Its lattice points have at most two parity types.
Consequently there are at most \(C_{\rm line}=4G\) nonempty cell/type
pieces on it, uniformly in \(n\), root and direction.
\end{lemma}
\begin{proof}
Inside the open square a nonaxis line crosses each of the
\(2(G-1)\) internal grid lines at most once. These crossings divide
it into at most \(2G-1\) intervals. Convexity prevents re-entry
into a cell. Coincident crossings at a vertex only reduce this count.
The integer points are \(p+kv\), with types alternating between
\(r\) and \(r+\bar v\). Thus there are at most
\(2(2G-1)\le4G\) pieces. Their index set is determined before the
random process is run.
\end{proof}

For an ordinary line \(\ell\ni p\) and one such piece \(E\), put
\[
 H_{\ell,E}(t)=r_\ell(t)
       \sum_{\substack{q\in(\ell\cap E)\setminus\{p\}\\q\text{ active}}}w_q(t),
 \qquad L_{\ell,E}=| (\ell\cap E)\setminus\{p\}|.
\]
Stop the fixed-label hierarchy at the first arrival at
\(H_{\ell,E}\ge L_{\ell,E}+en\), or the root boundary, whichever
comes first. As usual, observations at a boundary precede changes
of positive-retirement flags. For fixed \(m\), let
\(\Theta_m(n)=(1-\epsilon_n)^{-\binom{m+1}2}\).
Use the common-stop form \eqref{eq:mixedstop} with deterministic
horizon \(T\). Expansion by the number \(a\) of distinct leaves gives
\begin{equation}\label{eq:tagmoment}
 \E^{(p)}H_{\ell,E}^m
 \le\Theta_m(n)\sum_{a=1}^m
 \left\{\begin{matrix}m\\a\end{matrix}\right\}
 B^{m-a}L_{\ell,E}^a.
\end{equation}
The braces are Stirling numbers of the second kind. For distinct
leaves the exact connection is \(r^a\); for repetitions use
\(r^m\le r^a\), \(w^j\le B^{j-1}w\). There are no witnesses and
no designated-line connections, so no extra fixed multiplier is
needed. Ordered distinct-leaf tuples are counted by
\((L)_a\le L^a\).

\begin{lemma}[Finite height sum]\label{lem:height}
For \(e>0\), integers \(m\ge3\) and \(N\ge0\),
\begin{equation}\label{eq:height}
 \sum_{h=1}^N4h\left(\frac2{2+eh}\right)^m
 \le\frac{16}{e^2}\left(\frac2{2+e}\right)^{m-3}.
\end{equation}
\end{lemma}
\begin{proof}
Put \(a=e/2\). For \(h\ge1\),
\[
 h(1+ah)^{-m}\le\frac{(1+a)^{-(m-3)}}{a(1+ah)^2},
\]
and
\[
 (1+ah)^{-2}\le\frac1a
 \left[\frac1{1+a(h-1)}-\frac1{1+ah}\right].
\]
Sum the latter telescoping difference. Multiplication by four
gives \eqref{eq:height}, since \(4/a^2=16/e^2\).
\end{proof}

Apply Markov to \eqref{eq:tagmoment}, sum first over the at most
\(C_{\rm line}\) pieces of each line, then over directions of
height \(h\), using \(4h\) directions and \(L\le2n/h\).
The all-distinct term is bounded by Lemma~\ref{lem:height}.
Bounding the denominator in each repeated term by \((en)^m\)
gives the explicit remainder
\begin{align}\label{eq:tagremainder}
 R_m(n;e)=4e^{-m}\bigg[&2\left\{\begin{matrix}m\\1\end{matrix}\right\}
                 n(B/n)^{m-1}
 +4\left\{\begin{matrix}m\\2\end{matrix}\right\}
                 (1+\log n)(B/n)^{m-2}\notag\\
 &+\sum_{a=3}^{m-1}2^a
 \left\{\begin{matrix}m\\a\end{matrix}\right\}
 \left(1+\frac1{a-2}\right)(B/n)^{m-a}\bigg].
\end{align}
Thus the probability of any such first crossing, under the fixed
root's point law, is at most
\begin{equation}\label{eq:tagprob}
 C_{\rm line}\Theta_m(n)
 \left[\frac{16}{e^2}\left(\frac2{2+e}\right)^{m-3}
       +R_m(n;e)\right].
\end{equation}
No union over times, or over all roots, has been used.
For sufficiently large fixed \(m\), \(R_m(n;e)\to0\); it is enough
that \(1-(1-d)(m-1)<0\) and \(m>2\).

On the complementary good event, sum the cell bounds after multiplying
by their deterministic mask values, then integrate \eqref{eq:tagflow}.
The self point costs a further \(TB/n\). The result is
\begin{equation}\label{eq:tagcapacity}
 1-r_\ell\le\text{deterministic ordinary-line capacity}
       +C_{\rm line}Te+TB/n+O(n^{-1}).
\end{equation}
Choose \(e=(1-\rho-\lambda)/(4TC_{\rm line})\) when \(T>0\).
The zero-time case has no births and is trivial. For large \(n\),
the last two terms are at most \((1-\rho-\lambda)/4\).
Using capacity at most \(\lambda\), the right side is then strictly
less than \(1-\rho\), precluding a tag. For any requested positive
error, choose \(m\) large enough that the first term of
\eqref{eq:tagprob} is below half that error, and then \(n\) large
enough for \(\Theta_m(n)\) and the remainder. This proves
\begin{equation}\label{eq:tagtail}
 \sup_p\Prob^{(p)}(\text{tag is an early saving cause})=o(1).
\end{equation}
This linewise count also applies to the designated-line clock
argument (an axis line has at most \(2G\) cell/type pieces).
It does not replace the full number of cells in the uniform
bounded-variation estimates of Section~\ref{sec:geometry}.

\begin{proof}[Completion of the analytic reduction]
The weight-cap estimate \eqref{eq:captail}, tag estimate
\eqref{eq:tagtail}, and global-stop estimate \eqref{eq:globaltail}
give \(\sup_p\Prob^{(p)}(E_p)=o(1)\). Their events may overlap;
a union bound suffices. Proposition~\ref{prop:loss} implies
\(\E|S_n|\ge(1-o(1))nA_n\), whereas the exact nonnegative loss
identity also gives \(\E|S_n|\le nA_n\). By Lemma~\ref{lem:chord},
\(A_n=A+O(n^{-1})\), proving Theorem~\ref{thm:bridge}.

For fixed \(c<A\) the expectation exceeds \(cn\) for all large
\(n\). If every possible legal output had cardinality smaller than
\(cn\), its expectation would be smaller too, a contradiction.
Equivalently, with \(k=\lceil cn\rceil\) and \(|S_n|\le2n\),
\[
 \Prob(|S_n|\ge k)\ge
 \frac{\E|S_n|-(k-1)}{2n-(k-1)}>0.
\]
This is an existence argument; it does not supply a practical
algorithm producing the target density on moderate grids.
\end{proof}




\section{The finite input and its geometric checker}\label{sec:certificates}

This section specifies the 481-phase input and proves the geometric
upper bounds computed by its checker. The checker evaluates whole
closed root cells and whole angular intervals. Its input is an integer
field table and a deletion-rank table; neither sampled risks nor
search priorities are acceptance premises. The concrete arithmetic
results and execution record are given in Section~\ref{sec:numerics}.

\subsection{A schedule defined by reverse deletion ranks}\label{sec:rankinput}
Index the four types by \(r=0,1,2,3\), with
\(\typ(p)=(p_x\bmod2)+2(p_y\bmod2)\). The file
\texttt{field\_units.txt} contains 1024 nonnegative integers
\(u_{r,a,b}\), in that order, with \(0\le a,b<16\).
Set \(G=256\) and write
\(C_{ij}=(i/G,(i+1)/G)\times(j/G,(j+1)/G)\). Define
\begin{equation}\label{eq:rankfield}
 W_r^{(0)}\big|_{C_{ij}}
 =\frac{249}{250\cdot10^8}
       u_{r,\lfloor i/16\rfloor,\lfloor j/16\rfloor}.
\end{equation}
All fields and masks are zero on every fine-grid boundary and outside
the square. Thus the coarse table defines the values, but the boundary
convention is that of the fine grid.

The file \texttt{full\_rank.bin} contains \(4G^2=262144\)
little-endian signed 32-bit integers \(\nu_{r,i,j}\), in type, first
coordinate, second coordinate order. A zero-valued cell has rank
\(-1\); every positive cell has a rank in \(\{0,\ldots,480\}\).
For \(0\le j\le481\), let
\[
 W_r^{(j)}=W_r^{(0)}\ind_{\{\nu\ge j\}},
 \qquad E_r^{(j)}=W_r^{(0)}\ind_{\{\nu=j\}}
       \quad(0\le j<481).
\]
The indicators here are type-dependent cell indicators.
Then \(W^{(j)}-W^{(j+1)}=E^{(j)}\), and \(W^{(481)}=0\).
For every nonempty rank class set
\(\delta_j=\max_{r,i,k}E_r^{(j)}|_{C_{ik}}>0\).
Use the phases in decreasing rank order, from 480 to 0. On the phase
of duration \(\delta_j\), use the mask
\begin{equation}\label{eq:rankmask}
 f_r=E_r^{(j)}/\delta_j.
\end{equation}
It has rational values in \([0,1]\). This defines the chronological
schedule completely, including its durations and masks.

\begin{lemma}[Reverse deletion and the common deadline]\label{lem:reverse}
The field at the end of the rank-\(j\) phase is \(W^{(j)}\).
Every positive typed cell is used once, and its deadline is the end
of its rank's phase. Consequently, a risk bound for \(W^{(j)}\)
on every rank-\(j\) root cell is a bound for that cell's actual
common-deadline risk. The terminal field is \(W^{(0)}\), and
\[
 A=\frac14\sum_r\int W_r^{(0)},\qquad
 T=\sum_{j=0}^{480}\delta_j.
\]
\end{lemma}
\begin{proof}
At the indicated endpoint the phases already run are exactly those
with ranks \(k\ge j\). Their exposures sum to
\(\sum_{k=j}^{480}E^{(k)}=W^{(j)}\). This equality holds for all
four fields at the same time. A positive typed cell belongs to just
one rank class, so its last use is precisely that phase. Finally the
classes partition all positive typed cells. There are no overlapping
accepted-region areas to count with multiplicity.
\end{proof}

All intermediate fields are componentwise at most \(W^{(0)}\).
Capacities and squared-chord risk are monotone for nonnegative fields.
In particular, a terminal-field capacity bound controls the entire
schedule. The rank correspondence itself is checked before any
geometric calculation: the program rejects an out-of-range rank, a
mismatch between positivity and rank, or an empty requested phase.

\subsection{Exact maxima on a closed root cell}
Let \(F\) be a nonnegative rectangular step field, extended by zero
outside the unit square. For a direction \((1,u)\), \(u>0\), put
\[
 C(v;u)=\int_{\R}F(t,v+ut)\dd t.
\]
For a closed root box \(B=[a_0,a_1]\times[b_0,b_1]\), its possible
intercepts form the interval
\(l(u)=b_0-ua_1\le v\le h(u)=b_1-ua_0\).
At each field vertex \((a,b)\), combine the signed corner
coefficients of its incident rectangles, with exterior values zero.
The coefficients for one rectangle of value \(f\), in lower-left,
lower-right, upper-left, upper-right order, are \(-f,+f,+f,-f\).

\begin{lemma}[Corner representation]\label{lem:corners}
For \(u>0\),
\begin{equation}\label{eq:corners}
 C(v;u)=\frac1u\sum_{a,b}k_{a,b}(v-b+ua)_+.
\end{equation}
\end{lemma}
\begin{proof}
For the indicator of \((a_0,a_1)\times(b_0,b_1)\), integrate
\(\ind_{v+ut>b_0}-\ind_{v+ut>b_1}\) over \(a_0<t<a_1\).
For the first indicator the result is
\([(v+ua_1-b_0)_+-(v+ua_0-b_0)_+]/u\), and similarly for the
second. Add the identities with the rectangle weights and combine
shared corners. A nonaxis line meets the zeroed grid lines in finitely
many points, which do not affect this integral.
\end{proof}

\begin{proposition}[Exact closed-box envelope]\label{prop:envelope}
For fixed \(u>0\), the maximum chord of all roots in \(B\) is
\begin{equation}\label{eq:envelope}
 U_B(u)=\max\bigl\{C(v;u):
 v\in\{l(u),h(u)\}\cup
       \{b-ua:l(u)\le b-ua\le h(u)\}\bigr\}.
\end{equation}
The second set ranges over the field vertices. In particular this
maximum is obtainable by exact rational arithmetic at rational \(u\).
\end{proposition}
\begin{proof}
By \eqref{eq:corners}, the chord is continuous and piecewise affine
in \(v\), with knots among \(b-ua\). The maximum on an interval is
at an endpoint or at an included knot. Every possible root intercept
belongs to that interval, and every intercept in it is attained by
some root of the box.
\end{proof}

The implementation sorts projected vertices, combines equal knots,
and accumulates the slope and then the value of the chord numerator.
A segment tree returns the largest included knot value; the two
interval endpoints are evaluated separately. The zero value and zero
slope beyond the last knot are checked. Taking a maximum over all
knots, rather than only the root interval, also gives a valid bound
simultaneously for every root in the square.

\subsection{A rigorous angular-bin bound by vertical dilation}
For \(\delta\ge0\), let
\[
 (D_\delta F)(t,y)=\sup_{|z-y|\le\delta}F(t,z).
\]
\begin{lemma}[Tube domination]\label{lem:tube}
For \((x_0,y_0)\in[0,1]^2\) and \(|u-u_0|\le\delta\),
\[
 c_{F,1}((x_0,y_0),u)
 \le c_{D_\delta F,1}((x_0,y_0),u_0).
\]
\end{lemma}
\begin{proof}
Only \(0\le t\le1\) contributes. On this interval,
\(|(u-u_0)(t-x_0)|\le\delta\). The point of the original line is
therefore among those in the vertical supremum at the corresponding
point of the central line. Integrate the resulting pointwise bound.
\end{proof}

Choose \(N=2048\), \(H=2N\), \(\delta=1/H\), and
\(u_k=(2k+1)/H\), \(0\le k<N\). Lemma~\ref{lem:tube} gives
\begin{equation}\label{eq:tube}
 \sup_{x\in B}\int_0^1 c_{F,1}(x,u)^2\dd u
 \le\frac1N\sum_{k=0}^{N-1}
       \left(\sup_{x\in B}c_{D_\delta F,1}(x,u_k)\right)^2.
\end{equation}
This is a bound on all angles in each bin, not a quadrature
approximation whose error has been omitted.

The dilation is itself a finite rational rectangular field. In each
column, let the original cell values be \(f_j\), with exterior values
zero. Since \(\delta<1/(2G)\), the central strip of cell \(j\)
has value \(f_j\), and the strip of width \(2\delta\) about its
boundary \(j/G\) has value \(\max(f_{j-1},f_j)\). These strips
give an explicit decomposition, including the two outer boundaries.
Values on strip endpoints do not affect the positive-slope chord.
The program forms the integer type-pair field before this dilation;
it does not round each type's field separately.

All dilated vertices have coordinates \((a/H,b/H)\) with integers
\(a,b\). At \(u=p/H\), their projected intercepts are \(Hb-pa\).
The integer sweep computes a numerator whose raw-field chord is
that numerator divided by \(pH\). With rounding scale
\(Q_c=10^6\), the true scaled chord in \eqref{eq:rankfield} is
bounded in \(1/Q_c\) units by
\[
 L=\left\lceil\frac{249\,\text{numerator}}{25000\,pH}\right\rceil.
\]
This uses \(Q_c\,249/(250\cdot10^8)=249/25000\).
The bin contributes the integer upper bound \(\lceil L^2/N\rceil\)
in units \(1/Q_c^2\). All divisions round upward.

\subsection{Complete Farey partitions and rational integral bounds}
If the tube estimate is too large to accept a phase, the checker
uses the same residual field and every root box of that phase in a
second calculation. Its partition is the list of all reduced fractions
in \([0,1]\) with denominator at most \(G=256\).

A projected-knot ordering or a knot's membership in the root intercept
interval can change only at a slope joining two field or root-box
vertices. Those vertices are on the \(1/G\) grid, so every such slope
in \([0,1]\) has reduced denominator at most \(G\). On an open
interval between consecutive fractions, the candidate set in
\eqref{eq:envelope} is fixed and every candidate has the form
\(a+b/u\). Thus \(U_B(u)\) is convex as a function of \(1/u\).
Its endpoint secant in that coordinate bounds it from above.
Switches between maximizing candidates require no extra partition.
At a positive partition endpoint, continuity of the chord and of the
root intercept interval ensures that each one-sided candidate limit
is bounded by the endpoint maximum used in the calculation.

Let the consecutive positive fractions be \(\alpha=a/b\) and
\(\beta=c/d\), so \(bc-ad=1\) and \(\beta-\alpha=1/(bd)\).
Let \(A,B\) be upward-rounded endpoint chord values in units
\(1/Q_c\). The reciprocal-coordinate secant through these rounded
values remains an upper bound for \(Q_c U_B\).

\begin{lemma}[Rational affine integral bound]\label{lem:affine}
An upper bound for \(Q_c^2\int_\alpha^\beta U_B(u)^2\dd u\) is
the smaller of \(\max(A^2,B^2)/(bd)\) and
\begin{equation}\label{eq:affine}
 \begin{cases}
 (A^2+AB+B^2)/(3bd),&B\le A,\\[2pt]
 (A'^2+A'B'+B'^2)/(3bd),&B>A,
 \end{cases}
\end{equation}
where
\[
 A'=A+\frac{(B-A)bc}{(ad+bc)^2},\qquad
 B'=B+\frac{(B-A)ad}{(ad+bc)^2}.
\]
\end{lemma}
\begin{proof}
Write the reciprocal-coordinate secant as \(V(u)=s+t/u\).
It is nonnegative and monotone on the interval. If \(B\le A\),
then \(t\ge0\), so the endpoint secant in \(u\) is an upper bound.
If \(B>A\), then \(t<0\) and \(V\) is concave; its tangent at
\((\alpha+\beta)/2\) is an upper bound. Substitution and
\(bc-ad=1\) give tangent endpoint values \(A',B'\).
The integral of the square of an affine function with endpoint
values \(x,y\) on an interval of length \(1/(bd)\) equals
\((x^2+xy+y^2)/(3bd)\). The constant endpoint-maximum bound is
valid by monotonicity. Take the smaller of these two valid bounds.
\end{proof}

On the first interval \((0,1/G)\), every active candidate is bounded
and of the form \(a+b/u\), hence \(b=0\). Its maximum is constant,
so the first positive endpoint suffices. The value at angle zero
does not affect the angular integral; it does not replace the axis
capacity check below.

For an ordinary vertex \((i/G,j/G)\) at angle \(p/q\), the
projection is \(qj-pi\). The integer numerator is normalized by
\(pG\), and the rounded scaled chord is
\(\lceil249\,\text{numerator}/(25000\,pG)\rceil\).
The tangent values in \eqref{eq:affine} are retained as exact rational
numbers. Only the final integral bound is rounded upward to an integer
in units \(1/Q_c^2\). No logarithmic numerical integration is used.

\subsection{Four direction families, integer ranges and acceptance}
For each of the six unordered type pairs, the checker transforms the
field and root box by each of
\[
 (x,y),\quad(x,1-y),\quad(y,x),\quad(y,1-x).
\]
These give the four direction families of \eqref{eq:risk} and preserve
the major-coordinate integration parameter. Type labels are retained.
No transpose symmetry of the fields or symmetry of the random process
is assumed. Each type receives three pairs and four families.

If their twelve integral bounds sum to the integer \(J\), then
\begin{equation}\label{eq:certRisk}
 \mathcal R_r(W^{(j)};x)\le\frac{19J}{750Q_c^2}
 \quad\hbox{throughout the checked root box}.
\end{equation}
The acceptance test is
\[
 J\le\left\lfloor\frac{9999\cdot750Q_c^2}{190000}\right\rfloor
      =39469736842105.
\]
It implies \(\mathcal R_r\le9999/10000<1\). Each output row is
identified by type and fine-cell coordinates. The replay driver
compares the complete set of rows, and its cardinality, with the rank
class to exclude missing or repeated cells. It independently compares
the row integers to the rational risk threshold. Passing only selected
cells or retaining an unsuccessful tube bound is insufficient.

For this fixed input, the largest single-type integer is 628311458
and the largest type-pair integer is 1181489833. Deletion and dilation
cannot increase the latter bound. Its physical chord, in \(1/Q_c\)
units and rounded upward, is less than 13000000. A corner coefficient
has absolute value at most four times the pair maximum, below
\(2^{33}\). Even at the largest supported angular parameter
\(N=65536\), projected coordinates are below \(2^{35}\); a coarse
sum-of-absolute-values bound puts sweep intermediates below \(2^{86}\).
Actual chord numerators are below \(2^{65}\) before multiplication
by 249. The values and arithmetic used for chord conversion and
squaring are therefore within signed 128-bit range.
The sum over all bins of one tube integral is at most
\(13000000^2+N\); twelve such sums, and their risk numerator after
multiplication by 19, are also far below that range.
The same bound holds for a Farey integral with the number of intervals
in place of \(N\), using the endpoint-maximum bound on each interval.

For the Farey tangent put \(D=(ad+bc)^2\). Since
\(ad,bc\le G^2\), its rational numerator uses
\[
 A_L=AD+(B-A)bc,\qquad B_L=BD+(B-A)ad,
\]
and denominator \(3bdD^2\). The bound
\(A_L,B_L\le13000000(D+G^2)\) puts
\(A_L^2+A_LB_L+B_L^2\) below \(2^{117}\), and the denominator
below \(2^{86}\). The final ceiling addition also fits. Ordinary
integer division and signed 128-bit intermediates suffice for these
fixed data. These estimates are part of the soundness argument;
runtime undefined-behavior checks provide an additional implementation
check. They are not a claim that arbitrary int64 field tables would
be safe without a new range bound.

\subsection{Capacity: axes, finite directions and the infinite tail}
Exchange the two coordinates, transforming the field and root as
well, when the second coordinate of the direction is larger in
absolute value. Reverse the direction's orientation if necessary.
It then has the form \(v=(h,b)\), with \(|b|\le h\). Put
\(q=b/h\); the horizontal-coordinate chord below, divided by
\(2h\), is exactly the capacity in \eqref{eq:capacity}.
For a nonaxis slope \(q\ne0\), the chord of a rectangular cell
\((A,B)\times(C,D)\) at intercept \(v\) is
\[
 \pos{\min\left(B,\max\left(\frac{C-v}{q},\frac{D-v}{q}\right)\right)
 -\max\left(A,\min\left(\frac{C-v}{q},\frac{D-v}{q}\right)\right)}.
\]
It is continuous and piecewise affine in \(v\), with knots among
\(C-qA,C-qB,D-qA,D-qB\). Hence the maximum of the summed chord
is attained at an intercept through a grid vertex. The independent
capacity calculation uses this direct rectangle-intersection formula,
not the corner-sweep routine.

Axis directions need a different rule. In an open horizontal strip
the chord is its weighted row sum, whereas on a zeroed horizontal
boundary it is zero. Check all open-strip values, and use column sums
for the vertical direction. Further boundary zeroing on the fine grid
cannot increase these capacities: it removes finitely many points
from a nonaxis chord and may zero an entire axis chord.

Enumerate all 72 primitive unoriented directions of heights one
through seven. A primitive direction has nonzero parity, so it pairs
the four types into two unordered pairs; there are 144 finite
direction/type-pair checks. The 16 directions of height at most three
are designated. For every remaining direction with \(h\ge8\),
\[
 L_{v,r}(W^{(0)};x)
 \le\frac{\|W_r^{(0)}+W_{r+\bar v}^{(0)}\|_\infty}{2h}
 \le\frac{\|W_r^{(0)}+W_{r+\bar v}^{(0)}\|_\infty}{16}.
\]
Checking the six pair maxima therefore completes the infinite tail.
All these calculations may use the coarse \(16\times16\) envelope.

\begin{proposition}[Finite geometric checking contract]\label{prop:contract}
Suppose the specified field and rank tables satisfy the stated format
and positivity correspondence, every rank class is covered by its
checked closed cells, and the exact tube or Farey calculation proves
the common bound \(U_R\le9999/10000\). Suppose the capacity
calculation covers all finite directions and the tail, with bounds
\(U_H<K\) and \(U_L<\lambda<1-\rho\). If \(9999/10000<d<1\),
then the schedule \eqref{eq:rankmask} has the geometric input
properties \eqref{eq:heavyinput}--\eqref{eq:riskinput}, with uniform
positive margins and the exact mass given by its finite field sum.
\end{proposition}
\begin{proof}
The capacity calculation covers every line intercept and direction,
and terminal-field domination transfers it to all times.
Proposition~\ref{prop:envelope}, Lemma~\ref{lem:tube} and
Lemma~\ref{lem:affine} justify the angular bounds on every checked
closed cell. Equation~\eqref{eq:certRisk} gives the risk bound, while
Lemma~\ref{lem:reverse} identifies the four tested fields with the
four actual fields at that root's deadline. The finite maxima leave
the positive margins \(K-U_H\), \(\lambda-U_L\) and
\(d-9999/10000\), independent of \(n\). The rank classes partition
all positive typed cells, giving the mass without overlap ambiguity.
Proposition~\ref{prop:pair} and \eqref{eq:discrete481} then provide
the uniform discrete comparison used by the stochastic argument.
\end{proof}

The executable \texttt{independent/check\_phase.cpp}, with
\texttt{farey\_independent.inc}, implements these two angular bounds.
Its explicit rectangles, standard sorting, segment tree, ordinary
division, wider intermediates and four separate coordinate transforms
provide a second implementation of the geometric conditions.
Both implementations still use the mathematical lemmas proved here.
Input hashes identify the data; they do not prove the inequalities.
For a complete reproduction the checker is rebuilt and every phase
is evaluated from the fixed input, followed by per-cell coverage and
integer comparisons. Reusing cached phase reports without binding them
to the current inputs and checker is not such a reproduction.

\section{The 481-phase application and a fixed error budget}\label{sec:numerics}

\subsection{One concrete finite input}
We use the field and reverse ranks defined in Section~\ref{sec:certificates},
with spatial grid \(G=256\), time scale \(249/250\), and 481 nonempty phases.
There are 242,176 positive type/cells; each is used exactly once. By
Lemma~\ref{lem:reverse}, their reverse order defines masks in \([0,1]\),
and the checked background field is the cumulative field at the end of
the corresponding forward phase. The four types use a common endpoint.
This construction replaces the earlier candidate input; its mass is a
direct sum over disjoint type/cells, with no unverified union-area correction.

The exact mass and duration are
\begin{equation}\label{eq:input481numbers}
 A=\frac{3842853135627}{2560000000000}=1.501114506104296875,
 \qquad T=\frac{30088384743729}{25000000000}.
\end{equation}
The field is already time-scaled in these formulas. No further factor
\((249/250)^2\) is applied to its squared-chord risk. The coefficient
of the twelve integrals is \(D_*/24=19/750\).

\begin{center}
\begin{tabular}{@{}ll@{}}
\toprule
Quantity & Exact value or choice\\
\midrule
Target \(c\) & \(1501/1000\)\\
Designated clock cap \(K\) & \(255/128\)\\
Ordinary capacity cap \(\lambda\) & \(4/5\)\\
Ordinary credit floor \(\rho\) & \(1/8\)\\
Continuum risk bound \(r_*\) & \(9999/10000\)\\
Weight exponent \(d\) & \(19999/20000\)\\
Virtual exponent \(\gamma\) & \(39999/40000\)\\
Geometric and correction exponents \(\xi,\zeta\) & \(10^{-6},10^{-6}\)\\
Finite-sieve coefficient \(D_*\) & \(76/125\)\\
Spatial grid; designated direction count & \(256;\ 16\)\\
Total early-retirement probability budget & \(10^{-5}\)\\
\bottomrule
\end{tabular}
\end{center}

The capacity bounds for this same field are
\begin{align*}
 U_H&=\frac{159305594247}{80000000000},
 &K-U_H&=\frac{69405753}{80000000000},\\
 U_L&=\frac{2535641192953}{3200000000000},
 &\lambda-U_L&=\frac{24358807047}{3200000000000},\\
 U_{\rm tail}&=\frac{294190968417}{400000000000},
 &1-\rho-\lambda&=\frac3{40}.
\end{align*}
All margins are strict. The finite calculation covers 144 direction/type-pair
rows, including both axes; six tail bounds cover all heights at least eight.

The maximum in the freshly recomputed complete angular-replay records is
\begin{equation}\label{eq:recordedrisk481}
 r_{\rm rec}=\frac{93740574637267}{93750000000000},\qquad
 r_*-r_{\rm rec}=\frac{50362733}{93750000000000}>0.
\end{equation}
It occurs at reverse rank 205, type 2, microcell \((5,238)\).
Every positive type/cell is covered by exactly one phase record.
The checking algorithms and their soundness are described in
Section~\ref{sec:certificates}; the execution status for this edition
is stated separately in Section~\ref{sec:status}.

The finite sieve consists of the 175 odd primes from 3 to 1049, with
\[
 D_P=\frac34\prod_p(1-p^{-2})<\frac{76}{125}.
\]
For the discrete pair sum set
\begin{equation}\label{eq:adiscrete}
 a=r_*+\frac1{320000}=\frac{319969}{320000}<d.
\end{equation}
Equation~\eqref{eq:discrete481} gives
\(\mathcal B_n^{\Gamma^c}\le a\log n\) once
\(\log n\ge3\cdot10^{18}\). This explicitly pays the additive
lattice error; its coefficient is not silently left equal to \(r_*\).
The exponent conditions hold because
\[
 1-d-2\xi-\zeta=\frac{47}{10^6}>0,\qquad
 \gamma-d=\frac1{40000}>0.
\]

\subsection{Three global causes with fixed orders}
All moment orders are chosen before taking \(n\) to infinity.
For the pair-correction and virtual-deletion guards use
\(m_{\rm guard}=2^{18}=262144\). Including the factor \(B=n^d\)
for transport to the root point law, the powers in
\eqref{eq:etatail} and \eqref{eq:virtualtail} are respectively
\begin{align*}
 4-m+2m\xi+d(m-1)+m\zeta+d&=-8.320768,\\
 3-\gamma m+d(m-1)+\xi+d&=-3.553599.
\end{align*}
Each of these two probabilities is therefore less than \(2\cdot10^{-6}\),
uniformly in the root, for sufficiently large \(n\). Their constants
are finite for this fixed order; no unspecified constant is assigned
an assumed numerical value.

For designated-line clocks, \(C_{\rm line}=4G=1024\) and
\[
 \varepsilon_H=\frac{K-U_H}{8C_{\rm line}T}
 =\frac{115676255}{1314581588376682496}.
\]
Equations~\eqref{eq:clockexplicit}--\eqref{eq:clockthreshold} give
the same budget \(2\cdot10^{-6}\). In particular,
\(\log n\ge10^7\) suffices for this clock estimate. This is a
threshold for one cause, not a value of the final \(n_0\).

\subsection{The weight cap with a fixed moment order}
By \eqref{eq:hazarddecomp}, for every fixed integer \(m\), uniformly
in the root,
\[
 I_p\le Y_n+R_n,\qquad
 \|Y_n\|_{L^m(\Prob^{(p)})}\le(r_*+o_m(1))\log n,
 \qquad \E^{(p)}R_n=o(\log n),
\]
where \(Y_n,R_n\ge0\). Put
\[
 h=\frac{d-r_*}{4}=\frac1{80000},\qquad
 b=\frac{d+r_*}{2}=\frac{39997}{40000},\qquad
 m_{\rm cap}=2^{22}=4194304.
\]
After increasing \(n\) by a root-independent threshold, we have
\begin{equation}\label{eq:fixedriskmoment}
 \|Y_n\|_{m_{\rm cap}}\le(r_*+h/2)\log n,\qquad
 \E^{(p)}R_n\le10^{-6}h\log n.
\end{equation}
For \(q=(r_*+h/2)/(b-h)=159985/159986\), Bernoulli's inequality
gives \(q^{-159985}\ge2\). Since \(26\cdot159985\le m_{\rm cap}\),
Markov's inequality yields
\begin{equation}\label{eq:fixedcapbudget}
 \Prob^{(p)}(I_p>b\log n)
 \le q^{m_{\rm cap}}+10^{-6}
 \le2^{-26}+10^{-6}<2\cdot10^{-6}.
\end{equation}
The two terms pay for \(Y_n>(b-h)\log n\) and \(R_n>h\log n\)
separately.

Both axes are designated directions, so an active root has
\(\sum_{\ell\ni p,\,\ell\in\Gamma} b_{N_\ell}\ge1\) in
Lemma~\ref{lem:product}. Its self-correction
contribution is nonpositive. Each designated factor is at most
\(2/(2-K)=256\), and there are 16 designated directions. Thus
\[
 \log w_p\le I_p+16\log256<I_p+128.
\]
Reaching \(w_p=n^d\) requires \(I_p>d\log n-128\).
If \(\log n>5120000\), this implies the event in
\eqref{eq:fixedcapbudget}, so the fourth cause fits its budget.

\subsection{Ordinary-line tags with a fixed moment order}
Use the line-cell allowance
\[
 e=\frac{1-\rho-\lambda}{4TC_{\rm line}}
 =\frac{9765625}{641885541199552},\qquad
 b_{\rm tag}=131458159,\qquad k_{\rm tag}=87,
\]
and fix \(m_{\rm tag}=b_{\rm tag}k_{\rm tag}+3=11436859836\).
Since \(b_{\rm tag}e/2\ge1\), Bernoulli's inequality implies
\[
 \left(\frac2{2+e}\right)^{m_{\rm tag}-3}\le2^{-87}.
\]
The principal term in \eqref{eq:tagprob} is bounded by
\begin{equation}\label{eq:fixedtagbudget}
 \frac{16C_{\rm line}}{e^2}\,2^{-87}
 =\frac{100590099609629337889425049}
 {219902325555200000000000000000000}
 <\frac1{2000000}.
\end{equation}
For this fixed order, \(m>2\) and \(1-(1-d)(m-1)<0\); hence
the remainder in \eqref{eq:tagremainder} tends to zero. Increase
\(n\) until \(\Theta_m(n)\le2\) and
\(C_{\rm line}R_m(n;e)\le1/2000000\). The tag probability is then
less than
\[
 2\left[\frac{16C_{\rm line}}{e^2}2^{-87}
          +\frac1{2000000}\right]<2\cdot10^{-6}.
\]
The credit identity and \eqref{eq:tagcapacity}, with this same
\(e\), ensure that a tag cannot hit its floor without a cell crossing.
This large moment order is one fixed integer; no expansion into
that many monomials is part of the numerical verification.

\subsection{The final rational comparison}
There are five causes, each with probability less than \(2\cdot10^{-6}\)
under the appropriate root point law. The maximum of their finitely
many root-uniform thresholds gives
\begin{equation}\label{eq:fixedfivebudget}
 \sup_p\Prob^{(p)}(E_p)<10^{-5}
\end{equation}
for every sufficiently large \(n\). The causes need not be independent
or disjoint. This is a bound for each root under its own point law,
not a simultaneous-success probability for all roots under one law.

\begin{corollary}[The numerical lower bound]\label{cor:target}
For every sufficiently large integer \(n\), there is a set
\(S\subseteq\{0,\ldots,n-1\}^2\) with no three distinct collinear
points and \(|S|\ge\lceil1501n/1000\rceil\).
\end{corollary}
\begin{proof}
The finite input satisfies Proposition~\ref{prop:contract}, with the
values above. The lattice comparison and positive margins supply
the hypotheses of the stopped process. Equation~\eqref{eq:massgrid}
gives \(A_n\ge A(1-512/n)^2\) for every integer \(n\ge512\).
By Proposition~\ref{prop:loss} and \eqref{eq:fixedfivebudget},
\(\E|S_n|/n\ge(1-10^{-5})A_n\). For \(n\ge2^{24}\), in
addition to the probability thresholds,
\begin{align}
 \frac{\E|S_n|}{n}-\frac{1501}{1000}
 &\ge(1-10^{-5})A(1-2^{-15})^2-\frac{1501}{1000}\notag\\
 &=\frac{2165079963272809251797}
 {274877906944000000000000000}>0.
 \label{eq:exactfinalcomparison}
\end{align}
Every output is legal by Proposition~\ref{prop:construction}.
Some output therefore has cardinality at least \(\lceil1501n/1000\rceil\).
All comparisons apply to every sufficiently large integer \(n\),
with no parity or divisibility restriction.
\end{proof}

A legal set has at most two points per row, so \(|S_n|\le2n\).
Equation~\eqref{eq:exactfinalcomparison} also gives
\[
 \Prob\{|S_n|\ge\lceil1501n/1000\rceil\}>\frac1{64000}
\]
past the same thresholds: on the complementary event \(|S_n|<1501n/1000\),
and the rational margin divided by \(2-1501/1000\) exceeds \(1/64000\).
The proof does not compute the final \(n_0\) or provide a practical
bound for small-grid simulations.

\section{Verification scope and provenance}\label{sec:status}

\subsection{Finite computation and mathematical proof}
The input consists of 1,024 nonnegative field integers and 262,144 ranks.
The latter assign all 242,176 positive type/cells to 481 phases.
Source, data and replay records are distributed under \path{verification/}.
The checker evaluates all four direction families directly, without
assuming reflection or type symmetry. The 481 freshly accepted records
use 385 tube bounds, 95 Farey bounds, and one global tube bound.

\noindent\textbf{Fresh execution for this edition.}
On 10 November 2026 the checker was rebuilt using GCC 13.3.0 on
Ubuntu under WSL, with four OpenMP threads and undefined-behavior
sanitization enabled. No cached phase result was reused.
All 481 phases and 242,176 positive type/cells passed; coverage checking
found no missing, duplicate or failing cell. The complete angular
replay took 1,355.4 seconds (about 22.6 minutes).
All 481 accepted CSV files are byte-identical to the incoming dossier's
records, and the maximum is exactly \eqref{eq:recordedrisk481}.
No sanitizer error was reported. Capacity and clock-constant checks
also passed. The separate arithmetic ledger passed 69 checks,
including actual field/rank reconstruction and agreement with this run.

The reproduction driver \path{verification/reproduce.py} requires a new
output directory and records the compiler, flags, input and source hashes,
executable hash, commands, exit codes and result hashes. Unsafe cached
resumption has been removed from the Python wrapper. The mathematical
C++ sources and fixed input are unchanged from the reviewed independent
implementation. Execution provenance and all per-phase results are
included in \path{verification/evidence/}; the machine-specific executable
is rebuilt from source rather than distributed.

The soundness arguments in Section~\ref{sec:certificates} connect finite
integer computations to geometric bounds. Hashes establish which files
were used; hashes alone do not prove those bounds. Likewise, exact
arithmetic certifies finite rational inequalities and field/rank
reconstruction, not the stochastic lemmas. Those lemmas are proved
in the preceding sections as natural-language mathematics.

The proof uses fixed orders and a fixed error budget. It needs neither
growing-order moment estimates nor an explicit value of the final \(n_0\).
The available partial Lean development concerns earlier parts of the
project and does not formally certify this manuscript or its 481-phase
application. No independent human review or complete formal verification
is claimed.

\subsection{Companion files}
The editable manuscript is \path{main.tex} and its section files.
\path{data/parameters.json} gives exact parameters and finite-sieve factors.
\path{scripts/check_arithmetic.py} reconstructs the input mass, phase
durations and supports from the field and rank files, then checks
capacity margins, discrete-risk slack, fixed-order witnesses and the
final rational budget. Its report is \path{evidence/arithmetic.json}.
Geometric replay, capacity checking and execution provenance have
separate reports in \path{verification/evidence/}; the arithmetic
script does not replace them. Source identifiers and hashes are
retained in the companion provenance records.

\subsection{AI use and licensing}
This work involved substantial use of OpenAI ChatGPT/Codex and AI reviews
supplied by the project maintainer. AI systems contributed mathematical
proposals and derivations, numerical programs, exact arithmetic checks,
Lean definitions and proofs, build orchestration, review, translation
and exposition. Separate AI contexts reviewed parts of the work. This
English edition was assembled and checked by AI assistance; separate
AI reviews are not independent human peer review.

The human maintainer chose the direction and requested experiments and
revisions as a hobby project. This disclosure does not attribute every
mathematical or programming step to human authorship. No named author,
affiliation or personal contact address is supplied. The repository's
hosting account is public. The project material and this manuscript
are distributed under the MIT license in \path{LICENSE}; retain its
notices when redistributing.

\begin{thebibliography}{9}
\bibitem{preparation}
Project preparation notes, \emph{Integrated proof draft preceding
Lean formalization}, 8--9 November 2026 (Japanese).
\href{https://github.com/agolkjaoi/no-three-in-line-lean/blob/282b636c160d049e3ef4d82f619fd44d98f8cf1d/docs/research-notes/preparation/PROOF_FOR_LEAN_JA.txt}{Fixed-version text}.
Historical source for the process and stopping arguments.
\bibitem{followup}
Project first followup, 9 November 2026 (Japanese).
\href{https://github.com/agolkjaoi/no-three-in-line-lean/blob/282b636c160d049e3ef4d82f619fd44d98f8cf1d/docs/research-notes/first-followup/FOLLOWUP_FOR_REVIEW_JA.txt}{Fixed-version text}.
Historical source for remaining exposure and the fixed loss budget.
\bibitem{second}
Project second followup, 9 November 2026 (Japanese).
\href{https://github.com/agolkjaoi/no-three-in-line-lean/blob/282b636c160d049e3ef4d82f619fd44d98f8cf1d/docs/research-notes/second-followup/SECOND_FOLLOWUP_FOR_REVIEW_JA.txt}{Fixed-version text}.
Historical source for the finite sieve and linewise cell count.
\bibitem{formal}
\emph{no-three-in-line-lean}, partial formalization snapshot, 10 November 2026.
\href{https://github.com/agolkjaoi/no-three-in-line-lean/tree/282b636c160d049e3ef4d82f619fd44d98f8cf1d}{Source snapshot}.
This is a partial formalization, not a formal certificate of the present
numerical theorem.
\bibitem{dossier}
Project 481-phase dossier and fixed-budget closure, 10 November 2026.
Finite input and executable verification sources are included in this
edition. The relevant arguments are developed in English here; the
historical Japanese supplements are not required for the mathematical
argument.
\end{thebibliography}

\end{document}
